\chapter{Fiber Information Geometry and Inference Histories}
\label{ch:relational-inference}

This chapter asks how the collective free energy of \cref{ch:elbo} selects an
inference history of the agents at $c_*$ and what duration the Fisher geometry of their
law fibers assigns to it.  A chosen metric makes the positive ray of the natural
gradient select an oriented orbit, whose Fisher length is an information clock
invariant under reparameterization (\cref{thm:hist-fisher-clock-invariance}).  A
parameter-independent Markov channel projects scores and loses Fisher information by an
exact conditional-covariance defect (\cref{thm:cg-fisher-contraction}), so a same-path
Markov image of a history, such as the law of a meta-agent, accumulates no more duration
(\cref{thm:hist-record-clock-contraction}), while a coarse orbit recomputed from its own
free energy compares with the fine one only through an oriented semiconjugacy, whose
rate cancels from every duration
(\cref{thm:hist-oriented-semiconjugacy,thm:hist-duration-relation}).  A three-agent
parity witness shows that a collective free energy and its Fisher tensor are not
functions of the agents' marginal laws.  \Cref{ch:two-timescales} uses the mobility and
clock results, \cref{ch:coarsegraining} the Fisher-loss theorem, and
\cref{ch:agent-network-rg,ch:full-graph-meta-agent-vfe} the coarse configuration map and
the record clocks.  Longer proofs are deferred to \cref{app:deferred-proofs}.

\paragraph{Scope.}
Law fibers and kernel fibers are distinct, and no manifold, dominating measure, affine
chart, cone, or Fisher metric is inferred from the words belief or model.  The Fisher
metric is chosen: no invariance axiom that characterizes it
(\cref{prop:ig-fisher-characterizations}) is imposed, and no uniqueness is asserted.  The
parity hyperedge (\cref{def:hist-parity-hyperedge}) is supplied generative data, not
derived from pairwise locality or a canonical agent inventory, and the promoted parity
coordinate (\cref{cor:hist-promoted-parity-rank}) is not thereby an agent or an
intervention variable.  No fine-to-coarse semiconjugacy is asserted for coarse maps that
discard directions on which the fine free energy depends.  Open are correlated lifts
beyond the parity witness (\cref{open:hist-parity-lift-scope}), global Fisher quotients
and canonical block data (\cref{open:hist-retained-fisher-quotient}), the quantitative
meta-agent clock ratio (\cref{open:hist-meta-clock-ratio}), and the slow-manifold limit
of the information-discarding case (\cref{open:obs-slow-manifold}).  \status{NOT-CLAIMED}

\paragraph{Notation.}
Tangent vectors are $Z,Z_1,Z_2,Z_3$, and $\mathbf s_Z$ is the score in the direction $Z$;
$\xi_{\mathcal M}$ is the fundamental vector field of $\xi$ in the Lie algebra
$\mathfrak g$ of $G$ on a $G$-manifold $\mathcal M$.  A parameter-independent Markov
kernel is $K$, with law map $K_\star p=pK$ and joint law
$p\otimes K(dy,d\bar y)=p(dy)K(y,d\bar y)$.  A history is $Q$ with curve parameter $t$,
$\mathsf a$ is a positive rate factor, $\mathcal I_6$ indexes the bits of the parity
witness and $\operatorname{par}$ is its parity character, $\mathcal N(\mu,\varsigma^2)$ is
the normal family, and upright letters $\mathrm c,\mathrm c_1,\mathrm c_2,\mathrm q$ are
constants local to one statement.

\section{Fiber statistical tensors}
\label{sec:pb-statistical-tensors}
\label{sec:ig-typing}
\label{sec:ig-notclaimed}

At $c_*$ a frame $u$ in the $G$-torsor $(\mathscr P_G)_{c_*}$ identifies each law fiber
$\mathcal B_x$ with $(\mathcal E_x)_{c_*}$ through $\iota_u:\beta\mapsto[u,\beta]$, and
$\iota_{u\cdot g}=\iota_u\circ\widehat\rho_x(g)$ by \eqref{eq:geo-quotient-convention},
so a tensor on $\mathcal B_x$ defines one on $(\mathcal E_x)_{c_*}$ exactly when it is
invariant under the represented action.

\hypothesisheading{Regular channel models}{hyp:pb-regular-models}
Each law fiber $\mathcal B_x$, $x\in\{b,m\}$, and each coarse fiber $\bar{\mathcal B}_x$
used below is a finite-dimensional smooth parametrized-measure model with smooth
represented action, as in \cref{hyp:geo-smooth-tier}, differentiable in quadratic mean
with a positive-definite Fisher form; every score direction is integrable to the third
power, the second- and third-order score moments depend smoothly on the point, and the
model is dominated enough to differentiate the divergences of
\cref{sec:pb-divergence-jets} under the integral through third order.  The hypothesis
excludes singular statistical strata, families whose support moves without this
differentiability, and law fibers given only by a chart without a sample-space action.

\definitionheading{Fiber statistical tensors and the model-parameter Fisher tensor}{def:pb-fiber-statistical-tensors}
With $\mathbf s_Z$ the score of $\mathcal B_x$ at $p$ in the direction
$Z\in T_p\mathcal B_x$, the Fisher and Amari--Chentsov tensors of the channel and, for
the generative parameter $\theta$, structural data $X$, tangent vectors $Z_1,Z_2$ of the
parameter space of $\theta$, and a probability law $\mathsf r(dX)$ on structural data,
the model-parameter Fisher tensor of the evidence $p_{\theta,V}(o\given X)$ and its
average are
\begin{align}
g^F_{x,p}(Z_1,Z_2)&=\E_p[\mathbf s_{Z_1}\mathbf s_{Z_2}],
\label{eq:pb-fiber-fisher}\\
T^F_{x,p}(Z_1,Z_2,Z_3)&=\E_p[\mathbf s_{Z_1}\mathbf s_{Z_2}\mathbf s_{Z_3}],
\label{eq:pb-fiber-amari}\\
\mathsf I_X(\theta)(Z_1,Z_2)&=\E_\theta\left[\partial_{Z_1}\log p_{\theta,V}(O\given X)
\partial_{Z_2}\log p_{\theta,V}(O\given X)\mathrel{\big|}X\right],
\label{eq:ig-conditional-model-fisher}\\
\mathsf I_{\mathsf r}(\theta)&=\int\mathsf I_X(\theta)\mathsf r(dX).
\label{eq:ig-averaged-model-fisher}
\end{align}
Each of $\mathsf I_X$ and $\mathsf I_{\mathsf r}$ is called a metric only where it is
nondegenerate.

On the latent sample space $\mathsf Y$ an interaction energy is a function, not a metric;
the law fibers and the configuration manifolds of \cref{sec:hist-vfe-orbit-selection}
carry $g^F$; the parameter space of $\theta$ carries $\mathsf I_X$ or
$\mathsf I_{\mathsf r}$, of rank depending on $X$ or $\mathsf r$, one available metric for
$\theta$ in the fast and slow coordinates of \cref{def:obs-fast-slow-blocks}, whose
model-channel recognition coordinates carry $g^F_m$.  Equal components of a quadratic
energy on $\mathsf Y$ and a Fisher metric, as in the Gaussian mean sector of
Supplement~G, are a coincidence of coordinates, not an identification.

\propositionheading{Descent of the statistical tensors}{prop:pb-statistical-tensor-descent}
Under \cref{hyp:pb-regular-models}, every $\widehat\rho_x(g)$ preserves $g^F_x$ and
$T^F_x$, so $(\iota_u^{-1})^*g^F_x$ and $(\iota_u^{-1})^*T^F_x$ are tensors on
$(\mathcal E_x)_{c_*}$ that do not depend on the frame $u$; they are denoted by the same
symbols.  The sample-space origin of the action cannot be replaced by invariance of the
metric alone: for the normal family, whose Fisher metric is
$(d\mu^2+2d\varsigma^2)/\varsigma^2$, the map induced by $z\mapsto(1+z)/(1-z)$ on
$z=\mu/\sqrt2+\mathrm i\varsigma$ preserves $g^F$, fixes $(\mu,\varsigma)=(0,1)$, and
does not preserve $T^F$.  \status{ESTABLISHED}

The proof is in \cref{sec:proofs-fiber-geometry}.

\propositionheading{Invariance characterizations of the Fisher metric}{prop:ig-fisher-characterizations}
On the open simplices of positive laws on all finite sets, the family of Fisher
metrics is, up to one positive factor, the unique family of continuous Riemannian
metrics, one on each simplex, invariant under the congruent Markov embeddings between
simplices \citep[Theorem~11.1]{Cencov1982}, restated in
\citet[Proposition~3.19(3)]{Ay2015}; the uniqueness is a property of the family,
since the embeddings relate simplices of different dimensions.  On the cone of positive finite measures on a
finite set, with coordinates $\mathsf x_i$ and total mass $|\mathsf x|=\sum_i\mathsf x_i$,
the positive-definite metrics invariant under the same maps are exactly those with
components $g_{\mathsf x}(\partial_i,\partial_j)=\mathrm f_2(|\mathsf x|)
+\delta_{ij}|\mathsf x|\mathrm f_1(|\mathsf x|)/\mathsf x_i$ for smooth functions with
$\mathrm f_1>0$ and $\mathrm f_1+\mathrm f_2>0$ \citep{Campbell1986}; on tangent vectors
of the simplex, where $\sum_iZ_i=0$, the $\mathrm f_2$ term vanishes and the metric is
$\mathrm f_1(1)$ times the Fisher metric.  For parametrized measure models on general
sample spaces, every local continuous quadratic form field invariant under sufficient
statistics is a combination of the Fisher form and the square of the differential of
total mass with coefficients that are functions of the total mass, and on statistical
models, where the total mass is one, the Fisher form is the unique weakly continuous such
field up to a constant \citep[Main Theorem~2.10 and Remark~3.23]{Ay2015}; the finite
case of this characterization is Campbell's \citep[Section~1]{Le2017}.  Uniqueness up
to a constant also holds, on separable metrizable sample spaces with their Borel
sigma-algebras, for strongly continuous forms that are monotone under statistics
\citep[Theorem~1.4]{Le2017}, for diffeomorphism-invariant smooth weak Riemannian metrics
on smooth positive probability densities on a closed manifold of dimension at least two
\citep[Theorem in the introduction]{BauerBruverisMichor2016}, whose invariant bilinear
forms on all smooth positive densities are the combinations of the Fisher--Rao form and
the product of the total masses of the two tangent densities with constant
coefficients \citep[Main Theorem]{BauerBruverisMichor2016}, and, on exponential
families and their derived families, for Riemannian metrics invariant under independent
and identically distributed extensions and canonical sufficient statistics
\citep[Theorem~1]{Dowty2018}.  \status{STANDARD}

\section{Divergence jets and the orbit Fisher form}
\label{sec:pb-divergence-jets}

\definitionheading{Regular invariant contrast and its jets}{hyp:pb-contrast}
A regular invariant contrast on a law fiber $\mathcal B$ is a function
$\mathscr D:\mathcal B\times\mathcal B\to\R$, $C^3$ near the diagonal, with
$\mathscr D\geq0$ and, near the diagonal, $\mathscr D(p,q)=0$ exactly when $p=q$,
invariant under the simultaneous represented action,
$\mathscr D(\widehat\rho(g)p,\widehat\rho(g)q)=\mathscr D(p,q)$, and with
positive-definite metric jet $g_{\mathscr D}$ of \eqref{eq:pb-divergence-metric-jet}.  With $(Z,0)$ and $(0,Z)$ the lifts of
$Z\in T_p\mathcal B$ to the two factors, extended as constant coordinate fields in one
chart, the jets of $\mathscr D$ are
\begin{align}
g_{\mathscr D}(Z_1,Z_2)
&=-[(Z_1,0)(0,Z_2)\mathscr D]_{(p,p)},
\label{eq:pb-divergence-metric-jet}\\
\Gamma_{\mathscr D}(Z_1,Z_2,Z_3)
&=-[(Z_1,0)(Z_2,0)(0,Z_3)\mathscr D]_{(p,p)},
\label{eq:pb-divergence-first-connection-jet}\\
\Gamma_{\mathscr D}^*(Z_1,Z_2,Z_3)
&=-[(0,Z_1)(0,Z_2)(Z_3,0)\mathscr D]_{(p,p)},
\label{eq:pb-divergence-dual-connection-jet}\\
T^F_{\mathscr D}
&=\Gamma_{\mathscr D}-\Gamma_{\mathscr D}^*.
\label{eq:pb-divergence-amari-jet}
\end{align}
The two third-order arrays are connection coefficients in the chart, their difference is
a symmetric cubic tensor. Invariance makes the metric and cubic jets define tensors,
and the connection jets define affine connections, on the associated fiber
independently of the frame.  The contrast $(p,q)\mapsto(p-q)^4$
on a one-dimensional fiber is positive off the diagonal with zero second jet, so it is
excluded.

\propositionheading{KL divergence jets}{prop:pb-kl-divergence-jets}
On a fiber satisfying \cref{hyp:pb-regular-models}, the relative entropy
$\KL(p\Vert q)$ is a regular invariant contrast, its metric jet is $g^F$, and its cubic
jet $T^F_{\KL}$ is the Amari--Chentsov tensor $T^F$.  \status{ESTABLISHED}

\corollaryheading{Orbit Fisher form}{cor:pb-orbit-fisher-form}
Assume \cref{hyp:pb-regular-models} in one channel.  For $q\in\mathcal B$ define
\begin{equation}
\mathfrak k_q(\xi,\eta)=g^F_q\bigl(\xi_{\mathcal B}(q),\eta_{\mathcal B}(q)\bigr),
\qquad
\mathfrak g_q=\{\xi\in\mathfrak g:\xi_{\mathcal B}(q)=0\}.
\label{eq:pb-orbit-fisher-form}
\end{equation}
Then $\mathfrak k_q$ is positive semidefinite, its radical is exactly $\mathfrak g_q$,
and $\mathfrak g_q$ is the Lie algebra of the stabilizer $G_q$ of $q$.  For every
$\xi\in\mathfrak g$, as $t\to0$,
\begin{equation}
\KL\left(q\Vert\widehat\rho(\exp t\xi)q\right)
=\frac{t^2}{2}\mathfrak k_q(\xi,\xi)+O(t^3),
\label{eq:pb-orbit-kl-expansion}
\end{equation}
and for $\xi\in\mathfrak g_q$ the left side vanishes for every $t$.  The same expansion
holds for every regular invariant contrast, with $g^F$ replaced by its metric jet
$g_{\mathscr D}$.  \status{ESTABLISHED}

The proofs of \cref{prop:pb-kl-divergence-jets,cor:pb-orbit-fisher-form} are in
\cref{sec:proofs-fiber-geometry}.

In the running example of \cref{sec:geo-running-example}, a loop holonomy $\exp t\xi$
acting on an agent's law gives the loop a transported divergence
$\tfrac{t^2}{2}\mathfrak k_q(\xi,\xi)+O(t^3)$ that vanishes to all orders exactly when
$\xi\in\mathfrak g_q$, the second-order form of law coherence up to holonomy stabilizers
(\cref{def:geo-law-coherence}).

With the link as a variable, the orbit Fisher form is the Hessian of the transported divergence at an
exact link, and the minimizers of that divergence are unions of double cosets of the stabilizers, hence
not isolated when a stabilizer has positive dimension (\cref{prop:link-law-level-energy}).

\section{Markov coarse morphisms and the vertical Fisher defect}
\label{sec:pb-fisher-defect}
\label{sec:cg-fisher-contraction}

\theoremheading{Score projection and Fisher loss}{thm:cg-fisher-contraction}
Let $\{P_\lambda\}$ be a family of laws on $\mathsf Y$, indexed by an open subset of
$\R^d$, that is differentiable in quadratic mean at $\lambda_0$ with centered
square-integrable score $\mathbf s_{\lambda_0}:\mathsf Y\to\R^d$, and let
$K:\mathsf Y\rightsquigarrow\bar{\mathsf Y}$ be a Markov kernel that does not depend on
$\lambda$ and satisfies $K(y,\bar{\mathsf Y})=1$ for every $y$.  Under the joint law
$P_{\lambda_0}\otimes K$, the coarse family $\{P_\lambda K\}$ is differentiable in
quadratic mean at $\lambda_0$ with score
\begin{equation}
\bar{\mathbf s}_{\lambda_0}(\bar Y)
=\E\left[\mathbf s_{\lambda_0}(Y)\mid\bar Y\right]
\quad P_{\lambda_0}K\text{-almost surely},
\label{eq:cg-score-projection}
\end{equation}
and its Fisher information obeys
\begin{equation}
\mathsf I_{\bar{\mathsf Y}}(\lambda_0)\preceq\mathsf I_{\mathsf Y}(\lambda_0),
\qquad
\mathsf I_{\mathsf Y}(\lambda_0)-\mathsf I_{\bar{\mathsf Y}}(\lambda_0)
=\E\operatorname{Cov}\left(\mathbf s_{\lambda_0}(Y)\mid\bar Y\right)\succeq0.
\label{eq:cg-fisher-loss}
\end{equation}
Equality in a direction $z\in\R^d$ holds exactly when $z^\top\mathbf s_{\lambda_0}(Y)$
is $\bar Y$-measurable.  No common dominating measure is assumed.
\status{ESTABLISHED}

The proof is in \cref{sec:proofs-fiber-geometry}.

\remarkheading{Fisher equality is not recovery}{rem:cg-fisher-equality-not-recovery}
Equality in \eqref{eq:cg-fisher-loss} at one parameter value is local score
sufficiency, not recovery of the experiment.  Let $Y=(Y_1,Y_2)$ have independent
Bernoulli coordinates with success probabilities $\tfrac12+\tfrac\lambda4$ and
$\tfrac12+\tfrac{\lambda^2}4$, and let $K$ discard $Y_2$.  At $\lambda=0$ the score of
$Y_2$ vanishes, so the fine and coarse Fisher informations both equal $\tfrac14$, yet no
parameter-independent kernel $K'$ satisfies $(P_\lambda K)K'=P_\lambda$ for every
$\lambda$ in a neighborhood of zero: applying any fixed $K'$ to the
Bernoulli law $P_\lambda K$ makes every output event probability affine in
$\lambda$, whereas the required output event $Y_2=1$ has probability
$\tfrac12+\tfrac{\lambda^2}4$.  \citet{KaganShepp2005} give an insufficient statistic that
preserves Fisher information at every parameter value, and
\citet[Theorem~7]{Pollard2013} explains such examples by the criterion of
\cref{thm:cg-fisher-contraction}: preservation at $\lambda$ holds exactly when the score
is almost surely a function of the statistic.  \status{ESTABLISHED}

At $c_*$ the theorem applies to maps between fibers, for instance from the joint fiber
of a block of agents to the fiber of their meta-agent.

\definitionheading{Markov coarse morphism and its vertical defects}{hyp:pb-markov-morphism}
Let $\bar{\mathcal B}_x$ be a second law fiber in channel $x$, with represented action
$\widehat{\bar\rho}_x$ of $G$, associated bundle $\bar{\mathcal E}_x$, and frame maps
$\bar\iota_u$.  A Markov coarse morphism in channel $x$ at $c_*$ is a map
$\mathscr K:(\mathcal E_x)_{c_*}\to(\bar{\mathcal E}_x)_{c_*}$ with
$\bar\iota_u^{-1}\circ\mathscr K\circ\iota_u=K_\star|_{\mathcal B_x}$ for one, and hence
every, frame $u$, where $K:\mathsf Y_x\rightsquigarrow\bar{\mathsf Y}_x$ is a Markov
kernel that does not depend on the law parameter, satisfies $K(y,\bar{\mathsf Y}_x)=1$
for every $y$, maps $\mathcal B_x$ smoothly into $\bar{\mathcal B}_x$, and is
equivariant, $K_\star\circ\widehat\rho_x(g)=\widehat{\bar\rho}_x(g)\circ K_\star$ for
every $g$; no common dominating measure is assumed.  The vertical Fisher defect and the
Amari defect of $\mathscr K$ are
\begin{align}
\Delta_F^{\mathscr K}
&=g^F_x-(T\mathscr K)^*\bar g^F_x,
\label{eq:pb-vertical-fisher-defect}\\
\Delta_T^{\mathscr K}
&=T^F_x-(T\mathscr K)^*\bar T^F_x.
\label{eq:pb-amari-defect}
\end{align}
Equivariance makes the frame representative independent of $u$, since
$\bar\iota_{u\cdot g}^{-1}\mathscr K\iota_{u\cdot g}
=\widehat{\bar\rho}_x(g)^{-1}K_\star\widehat\rho_x(g)$, and by
\cref{prop:pb-statistical-tensor-descent} both defects are tensors on
$(\mathcal E_x)_{c_*}$.

Restriction to a submodel, precomposition of an energy with a map, a kernel that depends
on the law parameter, a fitted coarse approximation, and a comparison across a mismatch
of generative and recognition parameters are Markov coarse morphisms, and fall under the
results below, only when one kernel that realizes them is exhibited, along the whole history when histories are
compared.  One kernel serves every law on one fixed-precision fiber of a Gaussian
family; Supplement~G shows that Gaussian Galerkin aggregation has this form on each
fixed-precision fiber and no such form across precisions.

\corollaryheading{Vertical Fisher defect}{thm:pb-pullback-fisher-defect}
For a Markov coarse morphism $\mathscr K$ and every tangent vector $Z$ at $p$, under the
joint law $p\otimes K$,
\begin{align}
&\Delta_F^{\mathscr K}\succeq0,
\label{eq:pb-fisher-defect-positive}\\
&\Delta_F^{\mathscr K}(Z,Z)
=\E\operatorname{Var}\left(\mathbf s_Z(Y)\mid\bar Y\right),
\label{eq:pb-fisher-defect-score-variance}
\end{align}
and $\Delta_F^{\mathscr K}(Z,Z)=0$ exactly when $\mathbf s_Z(Y)$ is $\bar Y$-measurable.
\status{ESTABLISHED}

\noindent\emph{Proof.}
In a frame $\mathscr K$ is $K_\star$.  For a smooth curve $p_t$ in $\mathcal B_x$ with
velocity $Z$ at $p$, \cref{thm:cg-fisher-contraction} gives the image curve
$K_\star p_t$ the score $\E[\mathbf s_Z(Y)\mid\bar Y]$, which by uniqueness of the
quadratic-mean score is the score of $\bar{\mathcal B}_x$ in the direction $TK_\star Z$,
with second moment $\bar g^F_x(TK_\star Z,TK_\star Z)$; the defect identity \eqref{eq:cg-fisher-loss} for this one-parameter family is
\eqref{eq:pb-fisher-defect-score-variance}, with the stated equality clause.  $\square$

\theoremheading{Vertical Fisher defect cocycle}{thm:pb-fisher-defect-cocycle}
Let $\mathscr K_{01}:(\mathcal E_0)_{c_*}\to(\mathcal E_1)_{c_*}$ and
$\mathscr K_{12}:(\mathcal E_1)_{c_*}\to(\mathcal E_2)_{c_*}$ be smooth maps between
fibers carrying Fisher tensors $g^F_0,g^F_1,g^F_2$.  Then, unconditionally,
\begin{equation}
\Delta_F^{\mathscr K_{12}\circ\mathscr K_{01}}
=\Delta_F^{\mathscr K_{01}}
+(T\mathscr K_{01})^*\Delta_F^{\mathscr K_{12}},
\label{eq:pb-fisher-defect-cocycle}
\end{equation}
and the Amari defects obey the same identity.  If both maps are Markov coarse morphisms,
both terms of \eqref{eq:pb-fisher-defect-cocycle} are positive semidefinite.  No nonzero
Amari defect is sign definite, since
$\Delta_T^{\mathscr K}(-Z,-Z,-Z)=-\Delta_T^{\mathscr K}(Z,Z,Z)$; Markov coarse-graining
therefore gives no monotonicity for $T^F$, and $\Delta_T^{\mathscr K}=0$ is an additional
property of $\mathscr K$.  \status{ESTABLISHED}

\noindent\emph{Proof.}
Contravariance of pullback gives
$\Delta_F^{\mathscr K_{12}\circ\mathscr K_{01}}
=g_0^F-(T\mathscr K_{01})^*g_1^F+(T\mathscr K_{01})^*[g_1^F-(T\mathscr K_{12})^*g_2^F]$,
which is \eqref{eq:pb-fisher-defect-cocycle}, and the same computation applies to $T^F$.
Positivity is \cref{thm:pb-pullback-fisher-defect} for each Markov arrow, and oddness is
homogeneity of degree three.  $\square$

\section{Configuration manifolds and the exact free-energy lift}
\label{sec:hist-vfe-orbit-selection}

\definitionheading{Configuration manifold at $c_*$}{def:hist-configuration-manifold}
The configuration manifold $\mathcal Q_i$ of agent $i$ at $c_*$ is the product
$(\mathcal E_b)_{c_*}\times(\mathcal E_m)_{c_*}$ of its belief and model law fibers,
identified with $\mathcal B_b\times\mathcal B_m$ through its frame $u_i$, or an embedded
submanifold of that product (a regular submodel), with points $Q_i=(q_i^b,q_i^m)$.  The
collective configuration manifold is $\mathcal Q_V=\prod_{i\in V}\mathcal Q_i$, extended
by further coordinate manifolds, such as open simplices of attention rows, when those
coordinates move.  A configuration manifold carries a metric $\mathsf G$ that must be
specified.  The channel-weighted metric of agent $i$ is
$\mathsf G_i=w_bg^F_b\oplus w_mg^F_m$ with positive weights, under which a $C^1$ curve
$Q_i(t)$ has squared speed
\begin{equation}
\lVert\dot Q_i\rVert_{\mathsf G_i}^2
=w_bg^F_b(\dot q_i^b,\dot q_i^b)+w_mg^F_m(\dot q_i^m,\dot q_i^m);
\label{eq:hist-pointwise-clock-speed}
\end{equation}
on $\mathcal Q_V$ the metric is either block diagonal, $\bigoplus_i\mathsf G_i$, or the
pullback of a joint Fisher metric along an exact lift (\cref{hyp:hist-exact-vfe-lift}).

The channel-weighted and block-diagonal metrics are frame independent by
\cref{prop:pb-statistical-tensor-descent}. The pullback along an exact lift is frame
independent when the lift is equivariant in the sense of \cref{hyp:hist-exact-vfe-lift},
since then $\widehat\rho(g)^*\iota_i^*\mathsf G_{\mathscr R_B}
=\iota_i^*\widehat\rho_{\mathscr R}(g)^*\mathsf G_{\mathscr R_B}=\iota_i^*\mathsf G_{\mathscr R_B}$,
and is otherwise frame relative: it depends on the choice of right inverse
(\cref{prop:hist-general-score-residual}), and equivariant correlated lifts under a
connected group on continuous fibers are \cref{open:hist-parity-lift-scope}. The
base-direction version of these metrics is in Supplement~B.

\theoremheading{Fisher radical and quotient of a pulled-back tensor}{thm:hist-retained-fisher-quotient}
Let $\operatorname{ret}:\mathcal Q\to\mathcal Q'$ be a smooth map of finite-dimensional
manifolds, $g$ a positive-semidefinite symmetric tensor on $\mathcal Q'$, and
$h=\operatorname{ret}^*g$.  Then
\begin{equation}
\operatorname{rad}h_Q
=(d\operatorname{ret}_Q)^{-1}\left(\operatorname{rad}g_{\operatorname{ret}(Q)}\right),
\label{eq:hist-cds-fisher-radical}
\end{equation}
so $\operatorname{rad}h_Q=\ker d\operatorname{ret}_Q$ exactly when
$\operatorname{im}d\operatorname{ret}_Q\cap\operatorname{rad}g_{\operatorname{ret}(Q)}
=\{0\}$.  Where $\operatorname{rad}h$ has constant rank and $h$ is $C^1$,
$\operatorname{rad}h$ is a $C^1$ subbundle and $h$ induces a positive-definite form on
the quotient bundle $T\mathcal Q/\operatorname{rad}h$.  A quotient manifold additionally
requires the null leaves to form a Hausdorff manifold leaf space; when
$\operatorname{rad}h=\ker d\operatorname{ret}$ and the connected leaves form such a space
$\mathcal Q/{\sim}$ with quotient projection a smooth submersion, $h$ is the pullback of a positive-definite tensor on
$\mathcal Q/{\sim}$.  For a positive $C^3$ family on a finite set, with
$\operatorname{ret}$ the law map into the open simplex and $g$ its Fisher metric, $h$ is
the Fisher tensor of the family, of class $C^2$, positive definite exactly where
$\operatorname{ret}$ is an immersion, and there $\operatorname{ret}$ is locally
injective and $h$ has a $C^1$ Levi-Civita connection.  For softmax logits $z\in\R^k$
with $p=\operatorname{softmax}(z)$, $h(z)=\operatorname{diag}(p)-pp^\top$ has radical
$\R\mathbf 1$ at every $z$, and for an objective that depends on $z$ only through $p$
the solutions $Z$ of $hZ=\nabla\Fenergy$ form one class $Z_0+\R\mathbf 1$, whose
Moore--Penrose representative $h^+\nabla\Fenergy$ is Euclidean-orthogonal to
$\mathbf 1$.  \status{ESTABLISHED}

The proof is in \cref{sec:proofs-fiber-geometry}.

\openproblemheading{Global Fisher quotients and canonical block data}{open:hist-retained-fisher-quotient}
\Cref{thm:hist-retained-fisher-quotient} gives a positive-definite quotient bundle on
every constant-rank stratum and a quotient manifold once a Hausdorff leaf space is
supplied, and \cref{prop:hist-law-fisher-agentization-nogo} shows that the law and its
Fisher tensor alone select no splitting into agent blocks.  Open are a criterion,
checkable on the configuration manifolds of this chapter, under which the null foliation
of a pulled-back Fisher tensor has a Hausdorff manifold leaf space, and block data on
that quotient determined by a stated typing of agents and channels.  Such a criterion
with a proof, for instance exhibiting the leaves as the orbits of a free and proper group
action, together with a construction of natural block data from the typing or a
counterexample showing that the typing does not determine it, settles the problem.

\hypothesisheading{Standing configuration tier}{hyp:hist-standing-configuration-tier}
\label{hyp:hist-regular-metric-domain}%
Wherever this chapter asserts a natural-gradient vector field or a Fisher duration on a
configuration manifold $\mathcal Q$, either $\mathcal Q$ is a finite-dimensional
configuration manifold as in \cref{def:hist-configuration-manifold}, a finite product of
regular statistical manifolds or an embedded submodel of one, with a $C^1$
positive-definite metric $\mathsf G$, or $\mathcal Q$ is the quotient of such a manifold
by the null leaves of a $C^1$ positive-semidefinite Fisher tensor of constant rank with a
Hausdorff manifold leaf space, as in \cref{thm:hist-retained-fisher-quotient}, and the
objective is constant on those leaves.  The objective $\Fenergy$ is $C^2$, and every
minimal regular exponential family with its Fisher metric is an instance.  The hypothesis
excludes spaces of sections over a base, which belong to Supplement~B, semidefinite
metrics without a constructed quotient, and configuration metrics called exact joint-law
Fisher metrics without the lift of \cref{hyp:hist-exact-vfe-lift}.

\definitionheading{Exact multi-agent free-energy lift}{hyp:hist-exact-vfe-lift}
Fix a block $B$ of latent coordinates, an outside recognition marginal
$\mathbb Q_{B^c}$, and a manifold $\mathscr R_B$ of admissible conditional recognition
kernels $r_B(dy_B\given y_{B^c})$.  An exact lift of the configuration manifold
$\mathcal Q_i$ consists of a smooth extraction map
$\operatorname{ext}_i:\mathscr R_B\to\mathcal Q_i$ and a smooth right inverse
$\iota_i:\mathcal Q_i\to\mathscr R_B$, $Q_i\mapsto r_B^{Q_i}$,
$\operatorname{ext}_i\circ\iota_i=\operatorname{id}_{\mathcal Q_i}$, such that, for every
$Q_i$, the support, finiteness, and integrability hypotheses of
\cref{thm:obs-local-multiagent-elbo,thm:obs-local-global-potential} hold for
$\mathbb Q_{B^c}r_B^{Q_i}$, the lifted family is $C^2$, and for every chart
$\operatorname{ch}$ of $\mathcal Q_i$ with coordinates $\lambda$ and every compact set
$\mathsf C$ in its image there is $\mathrm e_{\mathsf C}\in L^1(\mathbb Q_{B^c})$ with,
for jointly measurable versions of the derivatives,
\begin{equation}
\sup_{\lambda\in\mathsf C}
\bigl|\partial_\lambda^\alpha\Fenergy^{\rm ext}_{B,o}
(r_B^{\operatorname{ch}^{-1}(\lambda)};y_{B^c})\bigr|
\leq\mathrm e_{\mathsf C}(y_{B^c})
\quad\text{for }\mathbb Q_{B^c}\text{-almost every }y_{B^c}
\text{ and }|\alpha|\leq2.
\label{eq:hist-chartwise-vfe-envelope}
\end{equation}
The lifted objective is the outside-averaged local free energy
\begin{equation}
\overline{\Fenergy}_{B,o}(Q_i)
=\E_{\mathbb Q_{B^c}}\left[\Fenergy^{\rm ext}_{B,o}\left(r_B^{Q_i};Y_{B^c}\right)\right],
\label{eq:hist-outside-averaged-local-vfe}
\end{equation}
a history of it is an exact free-energy history, and the exact recognition Fisher metric
of $\mathcal Q_i$ is $\mathsf G_i=\iota_i^*\mathsf G_{\mathscr R_B}$ on a nondegenerate
stratum, or its quotient as in \cref{thm:hist-retained-fisher-quotient}. The lift is
equivariant when $\mathscr R_B$ carries a declared represented action
$\widehat\rho_{\mathscr R}$ of the structure group, per agent or diagonal and compatible
with the frames of the outside marginal $\mathbb Q_{B^c}$, that preserves
$\mathsf G_{\mathscr R_B}$ and satisfies
$\operatorname{ext}_i\circ\widehat\rho_{\mathscr R}(g)=\widehat\rho(g)\circ\operatorname{ext}_i$
and $\iota_i\circ\widehat\rho(g)=\widehat\rho_{\mathscr R}(g)\circ\iota_i$; without such
an action, $\mathsf G_i$ is defined relative to the frames in which $\iota_i$ is written.

The lift carries every correlation that the configuration does not determine, and
$\operatorname{ext}_i\circ\iota_i=\operatorname{id}$ makes it a lift rather than a family
of laws indexed by configurations.  With the outside marginal held fixed,
\eqref{eq:hist-outside-averaged-local-vfe} differs from the collective free energy along
$Q_i\mapsto\mathbb Q_{B^c}r_B^{Q_i}$ by a term independent of $Q_i$
(\cref{thm:obs-local-global-potential}), so the two have the same differential and, for
one metric, the same natural-gradient ray.

\propositionheading{Joint-versus-weighted-marginal score residual and lift dependence}{prop:hist-general-score-residual}
Let $P_\lambda$ be a regular dominated family of joint laws of blocks
$(Y_1,\dots,Y_n)$ with square-integrable joint score
$\mathbf s=\nabla_\lambda\log p_\lambda(Y)$, embedded marginal scores
$\mathbf s_i=\E[\mathbf s\mid Y_i]$, and scalar weights $w_i$, and define
\begin{align}
\mathsf G_{\mathrm J}&=\E[\mathbf s\mathbf s^\top],
&\mathsf G_i&=\E[\mathbf s_i\mathbf s_i^\top],
&\mathsf G_w&=\sum_iw_i\mathsf G_i,
\label{eq:hist-score-metrics}\\
\mathbf r&=\mathbf s-\sum_i\mathbf s_i,
&\mathfrak K&=\E[\mathbf r\mathbf r^\top],
&\mathfrak C&=\sum_{i\ne j}\E[\mathbf s_i\mathbf s_j^\top].
\label{eq:hist-score-residual-tensors}
\end{align}
The cross tensor $\mathfrak C$ is symmetric and can be indefinite,
\begin{equation}
\mathsf G_{\mathrm J}-\mathsf G_w
=\mathfrak K-\mathfrak C+\sum_i(1-w_i)\mathsf G_i,
\label{eq:hist-joint-weighted-marginal-residual}
\end{equation}
and on a tangent subspace $\mathcal T$, $\mathsf G_{\mathrm J}=\mathsf G_w$ holds exactly
when the right side of \eqref{eq:hist-joint-weighted-marginal-residual} vanishes on
$\mathcal T\times\mathcal T$.  For two blocks and unit weights, in a direction $Z$ with
scores $\mathbf s_Z=Z^\top\mathbf s$ and $\mathbf s_{i,Z}=Z^\top\mathbf s_i$,
\begin{equation}
\lVert\mathbf s_Z\rVert^2
-\lVert\mathbf s_{1,Z}\rVert^2-\lVert\mathbf s_{2,Z}\rVert^2
=\lVert\mathbf s_Z-\mathbf s_{1,Z}-\mathbf s_{2,Z}\rVert^2
-2\langle\mathbf s_{1,Z},\mathbf s_{2,Z}\rangle,
\label{eq:hist-joint-versus-product}
\end{equation}
whose first term is nonnegative and whose second is signed.  If the two blocks are
independent under every $P_\lambda$, then $\mathsf G_{\mathrm J}=\mathsf G_1+\mathsf G_2$,
and $\mathsf G_{\mathrm J}=w_1\mathsf G_1+w_2\mathsf G_2$ forces $w_1=1$ as soon as some
direction has $\mathsf G_1(Z,Z)>0=\mathsf G_2(Z,Z)$, and symmetrically for $w_2$.  No
Loewner order between $\mathsf G_{\mathrm J}$ and $\mathsf G_1+\mathsf G_2$ holds in
general: for the bivariate normal location family with precision
$\begin{psmallmatrix}1&\mathrm c\\\mathrm c&1\end{psmallmatrix}$ and $0<\mathrm c<1$, the
difference is
$\begin{psmallmatrix}\mathrm c^2&\mathrm c\\\mathrm c&\mathrm c^2\end{psmallmatrix}$, with
value $2\mathrm c(1+\mathrm c)>0$ along $(1,1)$ and $2\mathrm c(\mathrm c-1)<0$ along
$(1,-1)$.  Distinct right inverses of one extraction map can induce distinct pullback
metrics: for $\mathcal N((\lambda_1,\lambda_2),\mathrm I_2)$, with Fisher metric
$d\lambda_1^2+d\lambda_2^2$, and the extraction $(\lambda_1,\lambda_2)\mapsto\lambda_1$,
the lifts $\lambda_1\mapsto(\lambda_1,0)$ and $\lambda_1\mapsto(\lambda_1,1)$ induce
$d\lambda_1^2$, whereas $\lambda_1\mapsto(\lambda_1,\lambda_1)$ induces $2d\lambda_1^2$.
\status{ESTABLISHED}

The proof is in \cref{sec:proofs-fiber-geometry}.

\subsection{Canonicity boundary of the lift}
\label{sec:hist-canonical-dependence-selection}

A lift supplies the dependence that the marginal configurations do not determine.  The
following finite results show that naturality under local kernels forces the product
law and that admitting correlated refinements leaves no natural selector.

\definitionheading{Finite marginal selection problem}{def:hist-finite-marginal-selection}
For a finite nonempty typed list $\mathcal X=(\mathcal X_i)_{i=1}^n$ of finite sets, let
$\mathcal J_{\mathcal X}=\Delta(\prod_i\mathcal X_i)$,
$\mathcal J^1_{\mathcal X}=\prod_i\Delta(\mathcal X_i)$, and
$\operatorname{mar}_{\mathcal X}:\mathcal J_{\mathcal X}\to\mathcal J^1_{\mathcal X}$ the
map that extracts the one-coordinate marginals.  A selector is a map
$\operatorname{sel}_{\mathcal X}:\mathcal J^1_{\mathcal X}\to\mathcal J_{\mathcal X}$ with
$\operatorname{mar}_{\mathcal X}\circ\operatorname{sel}_{\mathcal X}=\operatorname{id}$.
A local morphism is a tensor product $K=\bigotimes_iK_i$ of finite Markov kernels of the
same arity, and a kernel $K:\prod_i\mathcal X_i\rightsquigarrow\prod_j\mathcal Y_j$ that
changes arity is marginal compatible when
$\operatorname{mar}_{\mathcal Y}K_\star=\bar K\operatorname{mar}_{\mathcal X}$ for some
map $\bar K:\mathcal J^1_{\mathcal X}\to\mathcal J^1_{\mathcal Y}$.  A family of
selectors is natural on a class of morphisms when
$K_\star\operatorname{sel}_{\mathcal X}=\operatorname{sel}_{\mathcal Y}\bar K$ for every
morphism in the class; the class is part of the data, and not every finite kernel is
admitted.

\theoremheading{Tensor-local naturality selects the product law}{thm:hist-local-markov-product-uniqueness}
\label{cor:hist-maximal-product-preserving-category}%
If a family of selectors is natural under every local morphism, then
$\operatorname{sel}_{\mathcal X}(\mu_1,\ldots,\mu_n)=\bigotimes_{i=1}^n\mu_i$.  Among
classes of marginal-compatible kernels that contain every local morphism, a natural
selector exists exactly when every admitted kernel sends product laws to the product of
their image marginals, and then it is the product selector; the product-preserving
marginal-compatible kernels therefore form the largest such class.  \status{ESTABLISHED}

\noindent\emph{Proof.}
For each $i$ let the kernel from a one-point set send its point to $\mu_i$.  Their
tensor product sends the unique law on the one-point product to $\bigotimes_i\mu_i$,
while the induced marginal map sends the unique marginal tuple to $(\mu_i)_i$, so
naturality forces the product selector.  With it forced, naturality for an added kernel
says exactly that it sends each product law to the product of its image marginals.
$\square$

\theoremheading{Admitted correlated refinements refute an absolute selector}{thm:hist-correlated-refinement-nogo}
For a source bit $y$ and target bits $(y',y'')$, admit the split kernels
$K^{\mathrm{split}}_{\mathrm c}(y',y''\mid y)=\mathbf 1_{\{y'=y\}}(1+\mathrm c)/2$ if
$y''=y$ and $\mathbf 1_{\{y'=y\}}(1-\mathrm c)/2$ if $y''\ne y$, for
$\mathrm c\in\{1/3,1/2\}$.  Both are marginal compatible, and no single-valued selector
is natural under both; the argument uses no local naturality.  In particular
marginalization is not injective on positive laws, so no selector returns every joint
law from its marginals.  \status{ESTABLISHED}

\noindent\emph{Proof.}
The two kernels push the fair source law to laws with fair marginals and atom multisets
$\{1/3,1/3,1/6,1/6\}$ and $\{3/8,3/8,1/8,1/8\}$.  Naturality would require the value of
the target selector at the fair marginal pair to equal both laws, which differ even after
relabeling, and the two positive laws share their marginals.  $\square$

Selection relative to a positive reference law is a different problem. A
finite-relative-entropy minimizer over a convex family of moment-constrained laws,
when it exists, is unique by strict convexity of relative entropy. Both that minimizer
and, for a deterministic coarse map, the reference conditional completion of a
prescribed coarse marginal (\cref{lem:prob-kl-chain-rule}) depend on the reference
law as additional data.

\section{A finite correlated lift: the parity witness}
\label{sec:hist-finite-parity-lift}

\definitionheading{Three-agent parity lift}{def:hist-three-agent-parity-lift}
Let $\mathcal I_6=\{1,2,3\}\times\{b,m\}$, $y=(y_{i,x})\in\{0,1\}^{\mathcal I_6}$, and
$\lambda=(\lambda_{i,x})\in(0,1)^{\mathcal I_6}$, where typed agent permutations preserve
the channel label, and set
\begin{align}
&p_0(y;\lambda)
=\prod_{(i,x)}\lambda_{i,x}^{y_{i,x}}(1-\lambda_{i,x})^{1-y_{i,x}},
\qquad
\mathrm D(\lambda)=\prod_{(i,x)}\lambda_{i,x}(1-\lambda_{i,x}),
\label{eq:hist-parity-product-and-amplitude}\\
&p_\varkappa(y;\lambda)=p_0(y;\lambda)+\varkappa\operatorname{par}(y)\mathrm D(\lambda),
\qquad
\operatorname{par}(y)=(-1)^{\sum_{(i,x)}y_{i,x}},
\qquad|\varkappa|<1.
\label{eq:hist-parity-lift-law}
\end{align}
The map $\iota_\varkappa:(0,1)^6\to\Delta^\circ_{63}$ sends $\lambda$ to
$p_\varkappa(\cdot;\lambda)$.  The parity amplitude $\varkappa$ is fixed data of the
lift, not an inferred coordinate, and its hyperedge correlation is
$\bar\varkappa=\varkappa/64$.

\theoremheading{Positive smooth lift with paired relabeling symmetry}{thm:hist-three-agent-parity-lift}
For every $|\varkappa|<1$, the law \eqref{eq:hist-parity-lift-law} is strictly positive
and normalized, and every proper marginal is the corresponding marginal of $p_0$.
Consequently $\iota_\varkappa$ is a smooth right inverse of the map that extracts the six
one-bit marginals, and it is an immersion.  For each marginal tuple $\lambda$, the
positive laws with those marginals form a $57$-dimensional set, of which varying
$\varkappa$ traces one line.  Let $H_{\mathrm{pair}}=(\mathbb Z_2)^3\rtimes S_3$ act by
complementing both bits of a selected agent and by typed agent permutations, and act on
$\lambda$ by the same complements and permutations.  Then
\begin{equation}
\iota_\varkappa(g\lambda)=g_\#\iota_\varkappa(\lambda)
\qquad\text{for every }g\in H_{\mathrm{pair}},
\label{eq:hist-parity-paired-equivariance}
\end{equation}
whereas the complement $F_{i,x}$ of one bit satisfies
\begin{equation}
(F_{i,x})_\#p_\varkappa(\cdot;\lambda)=p_{-\varkappa}(\cdot;F_{i,x}\lambda).
\label{eq:hist-parity-odd-pseudoscalar}
\end{equation}
Odd complements are therefore equivariant only if $\varkappa$ changes sign with them, or
if $\varkappa=0$.  \status{ESTABLISHED}

The proof is in \cref{sec:proofs-fiber-geometry}.

\corollaryheading{Agent-block local--global free-energy identity}{cor:hist-parity-agent-block-vfe}
Fix an agent block $B_i=\{(i,b),(i,m)\}$ and the four outside parameters
$\lambda_{B_i^c}$, and let $r_{i,\lambda_{B_i}}=p_\varkappa(y;\lambda)/p_{0,B_i^c}
(y_{B_i^c};\lambda_{B_i^c})$ be the conditional law of $y_{B_i}$ given $y_{B_i^c}$.  For
any interior block
parameters $\lambda_{B_i}$ and $\lambda'_{B_i}$ and any strictly positive posterior
$\boldsymbol\Pi_o$ on the Boolean cube, all relative entropies are finite, and for every
block tangent $Z$
\begin{align}
&\Fenergy_o^{\rm ext}\bigl(p_\varkappa(\cdot;\lambda'_{B_i},\lambda_{B_i^c})\bigr)
-\Fenergy_o^{\rm ext}\bigl(p_\varkappa(\cdot;\lambda_{B_i},\lambda_{B_i^c})\bigr)
\notag\\
&\qquad=\E_{p_{0,B_i^c}}\bigl[
\Fenergy_{B_i,o}^{\rm ext}(r_{i,\lambda'_{B_i}};Y_{B_i^c})
-\Fenergy_{B_i,o}^{\rm ext}(r_{i,\lambda_{B_i}};Y_{B_i^c})\bigr],
\label{eq:hist-parity-agent-local-global-vfe}\\
&d\Fenergy_o^{\rm ext}[Z]
=\E_{p_{0,B_i^c}}\bigl[d\Fenergy_{B_i,o}^{\rm ext}[Z]\bigr],
\label{eq:hist-parity-agent-local-global-differential}
\end{align}
with both differentials evaluated at the same lifted conditional law.  The
outside-averaged block free energy is thus the restriction of one collective free energy
up to a constant.  \status{ESTABLISHED}

\noindent\emph{Proof.}
By \cref{thm:hist-three-agent-parity-lift} the outside marginal is exactly
$p_{0,B_i^c}$, independent of $\lambda_{B_i}$, and strict positivity makes
$r_{i,\lambda_{B_i}}$ a positive conditional law.  The finite state space and the
positive posterior make every relative entropy finite, so
\eqref{eq:hist-parity-agent-local-global-vfe} is \cref{thm:obs-local-global-potential}
with this outside marginal, and differentiating the finite smooth identity gives
\eqref{eq:hist-parity-agent-local-global-differential}.  $\square$

\corollaryheading{Parity-family Fisher residual}{cor:hist-parity-fisher-residual}
For the family of \cref{def:hist-three-agent-parity-lift}, every pair of bits is
independent, so $\mathfrak C=0$, and with unit weights
$\mathsf G_{\mathrm J}-\sum_{(i,x)\in\mathcal I_6}\mathsf G_{i,x}=\mathfrak K\succ0$ when
$\varkappa\ne0$.  At $\lambda=\tfrac12\mathbf 1$, and there with $\varkappa=\tfrac12$,
\begin{align}
\mathsf G_{\mathrm J}=\frac{4}{1-\bar\varkappa^2}\mathrm I_6,
\qquad
\mathsf G_{\mathrm J}-\sum_{(i,x)\in\mathcal I_6}\mathsf G_{i,x}
&=\frac{4\bar\varkappa^2}{1-\bar\varkappa^2}\mathrm I_6,
\label{eq:hist-parity-center-fisher}\\
\mathsf G_{\mathrm J}=\frac{65536}{16383}\mathrm I_6,
\qquad
\mathsf G_{\mathrm J}-\sum_{(i,x)\in\mathcal I_6}\mathsf G_{i,x}
&=\frac{4}{16383}\mathrm I_6,
\label{eq:hist-parity-half-fisher-residual}
\end{align}
and at the center a weighted marginal metric equals the joint one on the full tangent
space exactly when
\begin{equation}
w_{i,x}=\frac{1}{1-\bar\varkappa^2}
\qquad\text{for every }(i,x)\in\mathcal I_6,
\label{eq:hist-parity-center-weight-equality}
\end{equation}
so unit weights are exact only at $\varkappa=0$.  Since the one-bit marginal families do not depend on $\varkappa$, no
rule that reads only them returns every joint Fisher tensor.  \status{ESTABLISHED}

The proof is in \cref{sec:proofs-fiber-geometry}.

\corollaryheading{Promoted parity has one hidden direction}{cor:hist-promoted-parity-rank}
For the seven-parameter family $(\lambda,\varkappa)\mapsto p_\varkappa(\cdot;\lambda)$ on
$(0,1)^6\times(-1,1)$, the derivative of the joint law has rank seven everywhere, while
the one-bit marginal map has rank six and
\begin{equation}
\ker d\left(\operatorname{mar}\circ p\right)=\operatorname{span}\{\partial_\varkappa\}.
\label{eq:hist-cds-promoted-parity-kernel}
\end{equation}
Thus $\varkappa$ is identifiable from the joint law and invisible to the one-bit
marginals.  \status{ESTABLISHED}

\noindent\emph{Proof.}
The one-bit success probabilities are exactly $\lambda$, so a vanishing joint derivative
has zero $\lambda$-component, and the remaining derivative is a multiple of
$\partial_\varkappa p_\varkappa=\operatorname{par}\mathrm D$, which is nowhere zero on the
open domain; hence the joint derivative is injective.  The same identity gives rank six
for the marginal map and \eqref{eq:hist-cds-promoted-parity-kernel}.  $\square$

At the center $\lambda=\tfrac12\mathbf 1$, where $\mathrm D=4^{-6}$, the lift
\eqref{eq:hist-parity-lift-law} is
\begin{equation}
p_\varkappa\bigl(y;\tfrac12\mathbf 1\bigr)
=\tfrac1{64}\left(1+\bar\varkappa\operatorname{par}(y)\right).
\label{eq:hist-parity-center-law}
\end{equation}

\definitionheading{Three-agent parity hyperedge}{def:hist-parity-hyperedge}
At the center \eqref{eq:hist-parity-center-law}, take the uniform latent prior
$\operatorname{Unif}(y)=1/64$ and the binary record kernel
\begin{equation}
K_{\mathrm{hyp}}(1\given y)=\tfrac12\left(1+\bar\varkappa\operatorname{par}(y)\right),
\qquad
K_{\mathrm{hyp}}(0\given y)=\tfrac12\left(1-\bar\varkappa\operatorname{par}(y)\right).
\label{eq:hist-parity-hyperedge-kernel}
\end{equation}

\propositionheading{Exact posterior and tight free energy for the hyperedge}{prop:hist-parity-hyperedge-vfe}
The record $o=1$ has evidence $\tfrac12$, and its posterior is exactly
$p_\varkappa(\cdot;\tfrac12\mathbf 1)$.  For the free energy
$\Fenergy_1(Q)=\KL(Q\Vert\operatorname{Unif})-\E_Q\log K_{\mathrm{hyp}}(1\given Y)$ the
posterior attains the tight value $\Fenergy_1=\log2$, the product law of its one-bit
marginals is $\operatorname{Unif}$, and
\begin{equation}
\Fenergy_1(\operatorname{Unif})-\log2=-\tfrac12\log(1-\bar\varkappa^2).
\label{eq:hist-parity-product-lift-excess}
\end{equation}
Consequently, when $\varkappa\ne0$, no function of the one-bit marginals of the
recognition law, the one-bit marginals of the posterior, and the evidence equals the free
energy on every recognition law.  \status{ESTABLISHED}

The proof is in \cref{sec:proofs-fiber-geometry}.

\openproblemheading{Scope beyond the finite parity witness}{open:hist-parity-lift-scope}
\Cref{thm:hist-three-agent-parity-lift} gives one smooth correlated right inverse of the
one-bit marginal map, equivariant under paired complements and typed permutations, and
\cref{thm:hist-correlated-refinement-nogo} refutes an absolute selector on the admitted
split kernels, so that question is closed.  Open are the existence, for a stated regular
family of agent laws and a stated class of morphisms that contains
non-product-preserving kernels, of a smooth right inverse natural under that class, and
the existence of correlated lifts equivariant under a connected structure group acting on
continuous law fibers.  A construction with proof for a stated family and class, or a
no-go theorem for it, settles each part.

\section{Natural-gradient orbits}
\label{sec:hist-natural-gradient-oriented-orbit}

\definitionheading{Oriented unparameterized history and natural-gradient history}{def:hist-oriented-history}
A representative of a history in a configuration manifold $\mathcal Q$ is a $C^1$ curve
$Q:\mathcal T\to\mathcal Q$ on an open interval with nowhere-vanishing velocity, and two
representatives define the same oriented history when
\begin{equation}
\widetilde Q=Q\circ\phi,
\qquad
\phi:\widetilde{\mathcal T}\longrightarrow\mathcal T
\text{ an orientation-preserving $C^1$ diffeomorphism}.
\label{eq:hist-oriented-reparameterization-equivalence}
\end{equation}
The resulting abstract oriented curve $\mathscr H$ comes with its traversal map
$\jmath:\mathscr H\to\mathcal Q$, which retains order and multiplicity; for an embedded
representative $\mathscr H$ is its oriented image.  Under
\cref{hyp:hist-standing-configuration-tier} the natural gradient of $\Fenergy$ is defined
by the first line below, the descent ray at a noncritical $Q$ is
$\mathcal R^-_{\Fenergy}(Q)=\{-\mathsf a\operatorname{grad}^{\mathsf G}\Fenergy(Q):\mathsf
a>0\}$, and a natural-gradient history is an oriented history whose representatives
satisfy the second:
\begin{align}
\mathsf G\bigl(\operatorname{grad}^{\mathsf G}\Fenergy,Z\bigr)
&=d\Fenergy[Z]
\quad\text{for every }Z\in T\mathcal Q,
\label{eq:hist-fisher-gradient-definition}\\
\dot Q(t)
&=-\mathsf a(t)\operatorname{grad}^{\mathsf G}\Fenergy(Q(t)),
\qquad \mathsf a(t)>0 .
\label{eq:hist-vfe-ray-representative}
\end{align}
When $\Fenergy$ is the lifted objective of \cref{hyp:hist-exact-vfe-lift}, such a history
is an exact free-energy history.

\propositionheading{Scalar mobility preserves the oriented orbit}{prop:hist-scalar-mobility-orbit}
On a noncritical domain, multiplying the natural-gradient field by a smooth positive
function changes only the parameterization of its orbits.  A positive-definite
anisotropic mobility can change the orbit: on the Euclidean plane, for
$\Fenergy=\tfrac12(z_1^2+z_2^2)$ and $\mathsf M=\operatorname{diag}(1,2)$, gradient
descent from $(z_1^0,z_2^0)$ with $z_1^0z_2^0\ne0$ follows the line $z_2=(z_2^0/z_1^0)z_1$, whereas
$\dot z=-\mathsf M\nabla\Fenergy$ follows
\begin{equation}
z_2=z_2^0\left(z_1/z_1^0\right)^2 .
\label{eq:hist-anisotropic-orbit-change}
\end{equation}
A natural-gradient history is not in general a geodesic: on the same Euclidean plane the
natural-gradient flow of $\Fenergy=\tfrac12(z_1^2+2z_2^2)$ is $\dot z_1=-z_1$,
$\dot z_2=-2z_2$, whose orbits are the curves \eqref{eq:hist-anisotropic-orbit-change},
while the geodesics are straight lines.  \status{ESTABLISHED}

\noindent\emph{Proof.}
If $\mathscr V=-\operatorname{grad}^{\mathsf G}\Fenergy$ and $\mathscr V'=\mathsf
a\mathscr V$ with $\mathsf a>0$, the integral curves have the same tangent lines and
orientation, and the change of parameter defined by a positive integral of
$\mathsf a^{-1}$ carries one equation into the other.  In the example, $\dot z_1=-z_1$
and $\dot z_2=-2z_2$ keep $z_2/z_1^2$ constant, which is
\eqref{eq:hist-anisotropic-orbit-change}, while ordinary gradient descent keeps
$z_2/z_1$ constant, and for $z_1^0z_2^0\ne0$ the images differ; the same equations are
the Euclidean natural-gradient flow of $\tfrac12(z_1^2+2z_2^2)$.  $\square$

A rate ratio between coordinate blocks, such as fast beliefs and slow models, is an
anisotropic mobility (the example is the ratio $1:2$), so the selected orbit can depend
on it; the rate-separated flow is in \cref{sec:obs-fast-slow}.

\propositionheading{Noncritical free-energy orbits are embedded}{prop:hist-vfe-orbit-embedded}
Let $\mathcal Q$ be finite-dimensional, $\Fenergy$ of class $C^2$, and $\mathsf M$ a
$C^1$ field of positive-definite symmetric forms on the cotangent spaces, for instance
$\mathsf M=\mathsf a\mathsf G^{-1}$ for a $C^1$ metric and a $C^1$ function
$\mathsf a>0$.  Let $\gamma$ be a maximal integral curve of
$\mathscr V=-\mathsf M(d\Fenergy)$ through a noncritical point.  Then $\gamma$ meets no
critical point, $\Fenergy\circ\gamma$ is strictly decreasing, and $\gamma$ is an injective
immersion that is a homeomorphism onto its image, an embedded $C^1$ curve on which
$\Fenergy$ is a strictly monotone continuous coordinate, independent of parameterization
and gauge invariant whenever $\Fenergy$ is.  In particular the multiplicity clause of
\cref{def:hist-oriented-history} is vacuous for noncritical histories, and for the
natural gradient the Fisher clock of \cref{def:hist-fisher-clock} satisfies
$d\mathfrak t=-d\Fenergy/\lVert\operatorname{grad}^{\mathsf G}\Fenergy\rVert_{\mathsf G}$
along the orbit.  \status{ESTABLISHED}

\noindent\emph{Proof.}
The field $\mathscr V$ is $C^1$, so its integral curves are unique, and its zeros are the
critical points of $\Fenergy$; since the constant curve at a zero is an integral curve, a
curve through a nonzero of $\mathscr V$ never reaches one.  Along $\gamma$,
$\tfrac{d}{dt}\Fenergy(\gamma)=-\mathsf M(d\Fenergy,d\Fenergy)<0$, which makes
$\Fenergy\circ\gamma$ strictly decreasing, $\gamma$ injective, and $\dot\gamma\ne0$.  On
the image the inverse of $\gamma$ is $Q\mapsto(\Fenergy\circ\gamma)^{-1}(\Fenergy(Q))$, a
composite of continuous maps, so $\gamma$ is an embedding.  The clock relation is
\eqref{eq:hist-vfe-dissipation-per-length}.  $\square$

\propositionheading{A semidefinite tensor cannot be silently inverted}{prop:hist-semidefinite-gradient-obstruction}
If the proposed metric on a configuration space is only semidefinite, the
natural-gradient equation can have no solution or several.  \status{ESTABLISHED}

\noindent\emph{Proof.}
On $\R^2$ let $\mathsf G=dz_1^2$.  For $\Fenergy=z_2$ the equation
$\mathsf G(\operatorname{grad}\Fenergy,\cdot)=dz_2$ has no solution because the left side
annihilates $\partial_{z_2}$, and for $\Fenergy=z_1$ every
$\partial_{z_1}+\mathrm c\partial_{z_2}$ solves it.  $\square$

A quotient constructed as in \cref{thm:hist-retained-fisher-quotient}, or a separate
regularization, must precede any use of an inverse or a pseudoinverse.

\remarkheading{Discrete histories}{rem:hist-discrete-histories}
In a chart where the metric has a constant coefficient matrix $\mathsf G$, the explicit
Euler step of \eqref{eq:hist-vfe-ray-representative} with $\mathsf a\equiv1$ and step
$\Delta t>0$ is $Q_{n+1}=Q_n-\Delta t\mathsf G^{-1}\nabla\Fenergy(Q_n)$.  If
$\nabla^2\Fenergy\preceq\mathrm c_{\mathsf G}\mathsf G$ on the segment $[Q_n,Q_{n+1}]$,
with $\mathrm c_{\mathsf G}>0$, then
\begin{equation}
\Fenergy(Q_{n+1})
\leq\Fenergy(Q_n)
-\Delta t\left(1-\frac{\mathrm c_{\mathsf G}\Delta t}{2}\right)
\left\|\operatorname{grad}^{\mathsf G}\Fenergy(Q_n)\right\|^2_{\mathsf G},
\label{eq:hist-discrete-descent-bound}
\end{equation}
so every step with $0<\Delta t<2/\mathrm c_{\mathsf G}$ decreases the objective off the
critical set, and the bound is sharp: for a quadratic objective whose Hessian has largest
generalized eigenvalue $\mathrm c_{\mathsf G}$ relative to $\mathsf G$, an iterate on that
eigendirection is multiplied by $1-\Delta t\mathrm c_{\mathsf G}$, so
$\Delta t>2/\mathrm c_{\mathsf G}$ increases the objective there.  If, in the same chart,
$\mathscr V=-\mathsf G^{-1}\nabla\Fenergy$ is Lipschitz with constant
$\mathrm c_{\mathrm L}>0$ on the closed $\delta$-neighborhood of an exact arc $Q([0,t_1])$
and $\lVert\ddot Q\rVert\leq\mathrm c_2$ on the arc, then
\begin{equation}
\lVert Q_n-Q(n\Delta t)\rVert
\leq\frac{\mathrm c_2\Delta t}{2\mathrm c_{\mathrm L}}\left(e^{\mathrm c_{\mathrm L}n\Delta t}-1\right),
\qquad n\Delta t\leq t_1,
\label{eq:hist-euler-global-error}
\end{equation}
for every $\Delta t$ small enough that the right side stays below $\delta$ at
$n\Delta t=t_1$.  Euler iterates follow the selected orbit to first order in $\Delta t$
on compact parameter intervals without lying on it.  \status{ESTABLISHED}

\noindent\emph{Proof.}
Taylor's formula with integral remainder along the segment gives
$\Fenergy(Q_{n+1})\leq\Fenergy(Q_n)+\nabla\Fenergy^\top\Delta Q+\tfrac12\mathrm
c_{\mathsf G}\Delta Q^\top\mathsf G\Delta Q$ with
$\Delta Q=-\Delta t\mathsf G^{-1}\nabla\Fenergy$, which is
\eqref{eq:hist-discrete-descent-bound}; sharpness is the stated eigendirection
computation.  The local truncation error is at most $\mathrm c_2\Delta t^2/2$, so
$e_n=\lVert Q_n-Q(n\Delta t)\rVert$ obeys
$e_{n+1}\leq(1+\Delta t\mathrm c_{\mathrm L})e_n+\mathrm c_2\Delta t^2/2$ while the iterates
stay in the neighborhood, the discrete Gr\"onwall inequality gives
\eqref{eq:hist-euler-global-error}, and induction on $n$ keeps the iterates in the
neighborhood.  $\square$

\section{Fisher duration}
\label{sec:hist-fisher-duration}
\label{sec:pb-curve-taxonomy}
\label{sec:hist-clock-tiers}

\definitionheading{Fisher length and information clock}{def:hist-fisher-clock}
Let $Q:\mathcal T\to\mathcal Q$ be a $C^1$ curve in a configuration manifold with metric
$\mathsf G$.  At $c_*$ its velocity lies in the fibers, $\operatorname{ver}(\dot Q)=\dot
Q\in T_Q\mathcal Q$, so no connection enters.  Its Fisher speed, its length on
$[t_0,t_1]$, its cumulative clock from a chosen origin $t_0$, and, for a curve $\Gamma$
in one law fiber with its Fisher metric, the vertical Fisher length are
\begin{align}
\dot{\mathfrak t}(t)
&=\sqrt{\mathsf G\bigl(\dot Q(t),\dot Q(t)\bigr)},
\label{eq:hist-fisher-speed}\\
\operatorname{len}_{\mathsf G}(Q;[t_0,t_1])
&=\int_{t_0}^{t_1}\dot{\mathfrak t}(t) dt,
\label{eq:hist-fisher-length}\\
\mathfrak t_{Q,t_0}(t)
&=\int_{t_0}^{t}\dot{\mathfrak t}(t') dt',
\label{eq:hist-cumulative-fisher-clock}\\
\operatorname{len}_{g^F}(\Gamma)
&=\int\sqrt{g^F\bigl(\operatorname{ver}(\dot\Gamma),\operatorname{ver}(\dot\Gamma)\bigr)} dt .
\label{eq:pb-vertical-fisher-length}
\end{align}
The belief clock and the model clock of agent $i$ are the Fisher lengths of
$t\mapsto q_i^b(t)$ and $t\mapsto q_i^m(t)$, combined by the channel-weighted speed
\eqref{eq:hist-pointwise-clock-speed}.  Along a history of $\mathcal Q_V$ with
block-diagonal metric $\bigoplus_i\mathsf G_i$, the clock $\mathfrak t_i$ of agent $i$ is
the $\mathsf G_i$-length of $t\mapsto Q_i(t)$.  The origin is part of the clock, and the
parameter $t$ does not enter the length.

\theoremheading{Reparameterization invariance and unit-speed form}{thm:hist-fisher-clock-invariance}
Let $\widetilde Q=Q\circ\phi$ with $\phi:\widetilde{\mathcal T}\to\mathcal T$ an
orientation-preserving $C^1$ diffeomorphism.  Corresponding subarcs have the same length,
and
\begin{equation}
\mathfrak t_{\widetilde Q,\widetilde t_0}(\widetilde t)
=\mathfrak t_{Q,t_0}(\phi(\widetilde t)),
\qquad
t_0=\phi(\widetilde t_0).
\label{eq:hist-clock-reparameterization-law}
\end{equation}
If $\dot{\mathfrak t}>0$, the clock is strictly increasing and gives the unit-speed
representative of the oriented history; on a subinterval where $\dot{\mathfrak t}$
vanishes, as on a critical or null segment, the clock is constant and is not a
coordinate.  If $Q$ satisfies \eqref{eq:hist-vfe-ray-representative}, then on its
noncritical part
\begin{align}
\frac{d\mathfrak t}{dt}
&=\mathsf a\left\|\operatorname{grad}^{\mathsf G}\Fenergy\right\|_{\mathsf G},
\label{eq:hist-vfe-clock-rate}\\
\frac{dQ}{d\mathfrak t}
&=-\frac{\operatorname{grad}^{\mathsf G}\Fenergy}
{\left\|\operatorname{grad}^{\mathsf G}\Fenergy\right\|_{\mathsf G}},
\label{eq:hist-vfe-unit-tangent}\\
\frac{d\Fenergy}{d\mathfrak t}
&=-\left\|\operatorname{grad}^{\mathsf G}\Fenergy\right\|_{\mathsf G}.
\label{eq:hist-vfe-dissipation-per-length}
\end{align}
If the clock is read with a second metric $\mathsf G'$, for instance the unweighted
Fisher metric along a flow whose blocks carry rate weights, then along any $C^1$ history
with velocity $Z\ne0$
\begin{equation}
\left|\frac{d\Fenergy}{d\mathfrak t'}\right|
=\frac{|d\Fenergy(Z)|}{\lVert Z\rVert_{\mathsf G'}}
\leq\left\|\operatorname{grad}^{\mathsf G'}\Fenergy\right\|_{\mathsf G'},
\label{eq:hist-clock-metric-dissipation-bound}
\end{equation}
with equality exactly when $d\Fenergy=0$ or $Z$ is parallel to
$\operatorname{grad}^{\mathsf G'}\Fenergy$.
Along a history of $\mathcal Q_V$ with block-diagonal metric the squared speeds add,
$\dot{\mathfrak t}^2=\sum_i\dot{\mathfrak t}_i^2$, and over any common parameter interval
$\max_i\mathfrak t_i\leq(\sum_i\mathfrak t_i^2)^{1/2}\leq\mathfrak t\leq\sum_i\mathfrak
t_i$.  \status{ESTABLISHED}

The proof is in \cref{sec:proofs-fiber-geometry}.

\remarkheading{Endpoints do not determine an update duration}{prop:hist-realized-length-versus-distance}
For an update $Q$ from $Q(t_0)=Q'$ to $Q(t_1)=Q''$ in one path component of
$\mathcal Q$, the distance $\operatorname{dist}_{\mathsf G}(Q',Q'')$, the infimum of
$\operatorname{len}_{\mathsf G}(\Gamma;[t_0,t_1])$ over curves $\Gamma$ with these
endpoints, is at most $\operatorname{len}_{\mathsf G}(Q;[t_0,t_1])$, because the realized
update is one such curve, and the inequality is strict for a curved path in the Euclidean
Fisher geometry of the two-parameter normal location family.  Duration is also distinct
from the amount of curve parameter: in the unit-variance normal location family the path
$\mu(t)=e^{-t}$ has length $\int_0^\infty|\dot\mu(t)|dt=\int_0^\infty e^{-t}dt=1$ over an
unbounded parameter interval.  \status{ESTABLISHED}

\propositionheading{Finite duration for analytic objectives}{prop:hist-analytic-finite-duration}
Let $\mathcal Q$ be finite-dimensional with a continuous positive-definite metric
$\mathsf G$, and let $\gamma:[0,\infty)\to\mathcal Q$ satisfy
$\dot\gamma=-\operatorname{grad}^{\mathsf G}\Fenergy$.  If $\gamma(t)$ has an
accumulation point $Q_*$ as $t\to\infty$ and $\Fenergy$ is real analytic in a chart
around $Q_*$, then $\gamma(t)\to Q_*$, $Q_*$ is critical, and
$\operatorname{len}_{\mathsf G}(\gamma;[0,\infty))$ is finite.  This holds for every
forward history with relatively compact image when $\mathcal Q$ has a real-analytic atlas
in whose charts $\Fenergy$ is real analytic.  Analyticity cannot be weakened to
smoothness: for the $C^\infty$ function of \citet[Section~2, Eq.~(8)]{absil2005convergence}
on the Euclidean plane, the gradient history from the initial data given there is
bounded, has the unit circle as its $\omega$-limit set, and has infinite duration.
\status{ESTABLISHED}

\propositionheading{Invariant free-energy histories are horizontal}{prop:hist-invariant-history-horizontal}
Let a Lie group $\mathbb G$ act smoothly on a finite-dimensional configuration manifold
$\mathcal Q$ with a positive-definite metric $\mathsf G$, write
$\xi_{\mathcal Q}(Q)=\tfrac{d}{dt}\big|_{t=0}\exp(t\xi)\cdot Q$ and
$V_Q=\{\xi_{\mathcal Q}(Q):\xi\in\operatorname{Lie}\mathbb G\}$, and let $\Fenergy$ be
$C^1$ and $\mathbb G$-invariant.  Then $\operatorname{grad}^{\mathsf G}\Fenergy(Q)$ is
$\mathsf G$-orthogonal to $V_Q$ at every $Q$, every natural-gradient history has velocity
orthogonal to the orbits, and its gauge-quotient speed
$\min_{\xi}\lVert\dot Q-\xi_{\mathcal Q}(Q)\rVert_{\mathsf G}$ equals its Fisher speed.
If $\mathsf G$ is $C^\infty$ and the action is free, proper, and isometric, the quotient
map $\pi:\mathcal Q\to\mathcal Q/\mathbb G$ is a Riemannian submersion for a unique
metric $\bar{\mathsf G}$ on the quotient, $\Fenergy=\bar\Fenergy\circ\pi$, and
$T\pi\operatorname{grad}^{\mathsf G}\Fenergy=\operatorname{grad}^{\bar{\mathsf
G}}\bar\Fenergy\circ\pi$; each natural-gradient history descends to a natural-gradient
history of $\bar\Fenergy$ with the same Fisher duration.  \status{ESTABLISHED}

The proofs of \cref{prop:hist-analytic-finite-duration,prop:hist-invariant-history-horizontal}
are in \cref{sec:proofs-fiber-geometry}.

\section{Orthogonal clocks}
\label{sec:hist-global-clock}

\theoremheading{Exact criterion for an orthogonal clock}{thm:hist-global-clock-exactness}
Let $\mathcal U\subseteq\mathcal Q$ be a connected open set without critical points, with
$\Fenergy$ of class $C^2$ and $\mathsf G$ of class $C^1$.  Define $\mathsf n$,
$\mathbf e_{\Fenergy}$, and $\bm\alpha$ by the first line below, and call a $C^1$
function $\mathfrak t^\perp:\mathcal U\to\R$ that satisfies the second an orthogonal
clock:
\begin{align}
&\mathsf n=\left\|\operatorname{grad}^{\mathsf G}\Fenergy\right\|_{\mathsf G},
\qquad
\mathbf e_{\Fenergy}=-\frac{\operatorname{grad}^{\mathsf G}\Fenergy}{\mathsf n},
\qquad
\bm\alpha=\mathbf e_{\Fenergy}^\flat=-\frac{d\Fenergy}{\mathsf n},
\label{eq:hist-normalized-vfe-one-form}\\
&d\mathfrak t^\perp=\bm\alpha.
\label{eq:hist-global-clock-potential}
\end{align}
An orthogonal clock exists if and only if
\begin{equation}
d\bm\alpha=0,
\qquad
\oint_\gamma\bm\alpha=0
\quad\text{for every piecewise $C^1$ closed curve }\gamma\subset\mathcal U,
\label{eq:hist-global-clock-period-criterion}
\end{equation}
and on a simply connected $\mathcal U$ closedness suffices.  The local obstruction is
\begin{equation}
d\bm\alpha=\frac{1}{\mathsf n^2} d\mathsf n\wedge d\Fenergy,
\label{eq:hist-normalized-form-curvature}
\end{equation}
so $\bm\alpha$ is closed exactly when $\mathsf n$ is locally a function of $\Fenergy$,
and an orthogonal clock is then constant on every connected piece of every level set of
$\Fenergy$ and locally a strictly decreasing function of $\Fenergy$.  Whenever
$\mathfrak t^\perp$ exists, every unit-speed natural-gradient history in $\mathcal U$
obeys $d\mathfrak t^\perp/d\mathfrak t=1$, and the level sets of $\mathfrak t^\perp$ are
$\mathsf G$-orthogonal to $\mathbf e_{\Fenergy}$.  For $\Fenergy=z_1z_2$ on the open first
quadrant with the Euclidean metric,
\begin{equation}
d\bm\alpha
=\frac{z_1^2-z_2^2}{(z_1^2+z_2^2)^{3/2}} dz_1\wedge dz_2,
\label{eq:hist-nonexact-clock-example}
\end{equation}
which vanishes identically on no open set, so no orthogonal clock exists on any open
subset of the quadrant.  \status{ESTABLISHED}

\propositionheading{Unit-speed clocks from a global cross-section}{prop:hist-cross-section-clock}
Let $\mathcal U$ be as in \cref{thm:hist-global-clock-exactness}, and let
$\mathcal S\subset\mathcal U$ be a $C^1$ hypersurface that every maximal integral curve
of $\mathbf e_{\Fenergy}$ in $\mathcal U$ meets exactly once and transversally.  Then the
signed Fisher arc length along each orbit, measured from its crossing with $\mathcal S$,
is a $C^1$ function $\mathfrak t^{\mathcal S}:\mathcal U\to\R$ with
$d\mathfrak t^{\mathcal S}(\mathbf e_{\Fenergy})=1$, so every unit-speed natural-gradient
history obeys $d\mathfrak t^{\mathcal S}/d\mathfrak t=1$.  A level set of $\Fenergy$ that
every orbit meets is such a cross-section, and the level sets of
$\mathfrak t^{\mathcal S}$ are $\mathsf G$-orthogonal to $\mathbf e_{\Fenergy}$ only if
$\bm\alpha$ is closed.  For $\Fenergy=z_1z_2$ on the open first quadrant with the
Euclidean metric, the hyperbola $\{z_1z_2=1\}$ is a global cross-section, so a unit-speed
clock exists on the whole quadrant although no orthogonal one exists on any open subset.
\status{ESTABLISHED}

The proofs of \cref{thm:hist-global-clock-exactness,prop:hist-cross-section-clock} are in
\cref{sec:proofs-fiber-geometry}.

\section{Record clocks}
\label{sec:hist-record-clocks}

\corollaryheading{Same-path record and meta-agent clock contraction}{thm:hist-record-clock-contraction}
Let $t\mapsto P_t$ be a $C^1$ history in a family of laws on $\mathsf Y$ that is
differentiable in quadratic mean along it, with directional score $\mathbf s_t$, and let
$K:\mathsf Y\rightsquigarrow\bar{\mathsf Y}$ be a Markov kernel that does not depend on
$t$ and satisfies $K(y,\bar{\mathsf Y})=1$ for every $y$.  Then the record history
$\bar P_t=P_tK$, the image of the same fine path, satisfies
\begin{align}
\bar{\mathbf s}_t(\bar Y)
&=\E_t\left[\mathbf s_t(Y)\mid\bar Y\right],
\label{eq:hist-record-score-projection}\\
\dot{\mathfrak t}_{P}^2-\dot{\mathfrak t}_{\bar P}^2
&=\E_t\operatorname{Var}_t\left(\mathbf s_t(Y)\mid\bar Y\right)
\geq0,
\label{eq:hist-record-conditional-variance-defect}
\end{align}
and on every common parameter interval
\begin{equation}
\operatorname{len}(\bar P)\leq\operatorname{len}(P),
\qquad
\mathfrak t_{\bar P}(t_1)-\mathfrak t_{\bar P}(t_0)
\leq
\mathfrak t_{P}(t_1)-\mathfrak t_{P}(t_0),
\label{eq:hist-record-clock-contraction}
\end{equation}
with equality, when the fine length on the interval is finite, exactly when the fine
score is $\bar Y$-measurable for almost every $t$ in the interval. Equality of two
infinite lengths gives no such conclusion. The inequality applies to a meta-agent exactly when its law is the image
of the same fine history under one such kernel.  Parameter independence of the kernel
cannot be dropped: for a dominated channel $K_t$ with density
$\mathsf k_t(\bar y\given y)$, assume the joint family $P_t(dy)K_t(y,d\bar y)$ is
differentiable in quadratic mean with square-integrable score
$\mathbf s_t(y)+\partial_t\log\mathsf k_t(\bar y\given y)$. Its record score is
\begin{equation}
\bar{\mathbf s}_t(\bar Y)
=\E_t\left[
\mathbf s_t(Y)
+\partial_t\log\mathsf k_t(\bar Y\given Y)
\mathrel{\big|}\bar Y
\right],
\label{eq:hist-parameter-dependent-channel-score}
\end{equation}
and a constant input followed by a channel that emits a Bernoulli bit of success
probability $(1+e^{-t})^{-1}$ has zero fine speed and record speed
$(2\cosh(t/2))^{-1}>0$.  \status{ESTABLISHED}

The proof is in \cref{sec:proofs-fiber-geometry}.

\begin{figure}[htbp]
\centering
\begin{tikzpicture}[>=stealth]
  \node[draw,rounded corners,fill=lightblue,minimum width=3.3cm,minimum height=0.8cm]
    (fine) at (0,2.1) {fine history $P_t$};
  \node[draw,rounded corners,fill=lightgreen,minimum width=3.8cm,minimum height=0.8cm]
    (image) at (6.2,2.1) {same-path image $P_tK$};
  \draw[->,very thick,primaryblue] (fine)--(image);
  \node[text=primaryblue,font=\small] at (3.1,2.85)
    {parameter-independent Markov channel: $d\mathfrak t_{PK}\leq d\mathfrak t_P$};
  \node[draw,rounded corners,fill=lightblue,minimum width=3.3cm,minimum height=0.8cm]
    (fflow) at (0,0) {$\operatorname{Fl}^{\mathscr V_\ell}_t(Q_0)$};
  \node[draw,rounded corners,fill=lightgreen,minimum width=3.8cm,minimum height=0.8cm]
    (mflow) at (6.2,0) {$\operatorname{Fl}^{\mathscr V_{\ell+1}}_{t'}(\mathsf F_\ell Q_0)$};
  \draw[->,very thick,dashed,sciencegreen] (fflow)--(mflow);
  \node[text=sciencegreen,font=\small] at (3.1,0.75)
    {independent flows require oriented semiconjugacy};
\end{tikzpicture}
\caption{The upper arrow is the same-path comparison of
\cref{thm:hist-record-clock-contraction}, with its conditional-variance defect.  The
lower arrow compares independently generated flows and exists only under the
semiconjugacy condition of \cref{def:hist-oriented-semiconjugacy}.}
\label{fig:hist-markov-clock-and-semiconjugacy}
\end{figure}

\openproblemheading{Quantitative meta-agent clock ratio}{open:hist-meta-clock-ratio}
\Cref{thm:hist-record-clock-contraction} shows that a meta-agent whose law is a Markov
image of its constituents' joint history accumulates no more Fisher duration than that
history, with the exact defect \eqref{eq:hist-record-conditional-variance-defect}.  Open
is the ratio between the meta-agent's duration and the constituents' collective duration
for an $N$-constituent meta-agent of \cref{ch:agent-network-rg}, as a function of $N$ and
of the correlation between the constituents' velocities, for coherent histories (aligned
velocities) and incoherent ones (independent random velocities).  The expected ratio of
arc lengths is not the square root of the expected ratio of relative entropies, so a
divergence computation does not settle it; a theorem that computes or bounds the ratio
for a stated coarse map and a stated law of constituent velocities, including its
dependence on $N$ in both regimes, does.

\section{The coarse configuration map and oriented semiconjugacy}
\label{sec:hist-configuration-coarse-map}
\label{sec:hist-independent-flows}

\definitionheading{Coarse configuration map}{def:hist-coarse-configuration-map}
A coarse configuration map at depth $\ell$ is a smooth map
\begin{equation}
\mathsf F_\ell:\mathcal Q_\ell\longrightarrow\mathcal Q_{\ell+1}
\label{eq:hist-configuration-coarse-map}
\end{equation}
from the configuration manifold of the fine agents at $c_*$ to that of their meta-agents,
specified separately from the sample kernels that may realize it.  The fine and coarse
natural-gradient fields are
$\mathscr V_\ell=-\operatorname{grad}^{\mathsf G_\ell}\Fenergy_\ell$ and
$\mathscr V_{\ell+1}=-\operatorname{grad}^{\mathsf G_{\ell+1}}\Fenergy_{\ell+1}$, with
local flows $\operatorname{Fl}^{\mathscr V_\ell}_t$ and
$\operatorname{Fl}^{\mathscr V_{\ell+1}}_{t'}$.

\definitionheading{Oriented semiconjugacy}{def:hist-oriented-semiconjugacy}
For an open $\mathcal U\subseteq\mathcal Q_\ell$, the pair $(\mathsf F_\ell,\mathsf
a_\ell)$ with $\mathsf a_\ell:\mathcal U\to(0,\infty)$ continuous is an oriented
semiconjugacy of $\mathscr V_\ell$ onto $\mathscr V_{\ell+1}$ over $\mathcal U$ when
\begin{equation}
T_Q\mathsf F_\ell \mathscr V_\ell(Q)
=\mathsf a_\ell(Q) \mathscr V_{\ell+1}(\mathsf F_\ell Q)
\quad\text{for every }Q\in\mathcal U.
\label{eq:hist-oriented-vector-field-semiconjugacy}
\end{equation}

\lemmaheading{The factor is determined and regular off the coarse critical set}{lem:hist-semiconjugacy-factor}
If $\mathscr V_{\ell+1}(\mathsf F_\ell Q)\ne0$, then $\mathsf a_\ell(Q)$ is unique and
\begin{equation}
\mathsf a_\ell(Q)
=\frac{\mathsf G_{\ell+1}
\left(T_Q\mathsf F_\ell \mathscr V_\ell(Q),\mathscr V_{\ell+1}(\mathsf F_\ell Q)\right)}
{\left\lVert\mathscr V_{\ell+1}(\mathsf F_\ell Q)\right\rVert^2_{\mathsf G_{\ell+1}}},
\label{eq:hist-semiconjugacy-factor-formula}
\end{equation}
so $\mathsf a_\ell$ is $C^k$ wherever $Q\mapsto T_Q\mathsf F_\ell\mathscr V_\ell(Q)$,
$Q\mapsto\mathscr V_{\ell+1}(\mathsf F_\ell Q)$, and the pulled-back coefficients of
$\mathsf G_{\ell+1}$ are $C^k$, for instance when $\mathsf F_\ell\in C^{k+1}$ and
$\mathscr V_\ell,\mathscr V_{\ell+1},\mathsf G_{\ell+1}$ are $C^k$; continuity of
$\mathsf a_\ell$ is therefore automatic where $\mathscr V_{\ell+1}\ne0$ on
$\mathsf F_\ell(\mathcal U)$.  If $\mathscr V_{\ell+1}(\mathsf F_\ell Q)=0$, then
\eqref{eq:hist-oriented-vector-field-semiconjugacy} forces
$T_Q\mathsf F_\ell\mathscr V_\ell(Q)=0$ and fixes no value of $\mathsf a_\ell(Q)$; limits
from the noncritical set may determine it or may obstruct any continuous positive
extension.  \status{ESTABLISHED}

\noindent\emph{Proof.}
Pair \eqref{eq:hist-oriented-vector-field-semiconjugacy} with
$\mathscr V_{\ell+1}(\mathsf F_\ell Q)$ and divide by its positive squared norm.
$\square$

\theoremheading{Oriented semiconjugacy and maximal intervals}{thm:hist-oriented-semiconjugacy}
\label{prop:hist-oriented-semiconjugacy}%
Assume \eqref{eq:hist-oriented-vector-field-semiconjugacy} on an open $\mathcal U$ with
$\mathscr V_{\ell+1}$ locally Lipschitz.  Fix $Q\in\mathcal U$, let $\mathcal T_Q$ be the
connected component of $0$ in the set of $t$ for which
$\operatorname{Fl}^{\mathscr V_\ell}_{t''}(Q)\in\mathcal U$ for every $t''$ between $0$
and $t$, and set $\psi_Q(t)=\int_0^t\mathsf a_\ell(\operatorname{Fl}^{\mathscr
V_\ell}_{t''}Q) dt''$.  Then $\psi_Q$ is a strictly increasing $C^1$ diffeomorphism of
$\mathcal T_Q$ onto an open interval $\mathcal T'_Q\ni0$, and
\begin{equation}
\mathcal T'_Q\subseteq\bar{\mathcal T}^{\max}_{\mathsf F_\ell Q},
\qquad
\mathsf F_\ell\left(\operatorname{Fl}^{\mathscr V_\ell}_t(Q)\right)
=\operatorname{Fl}^{\mathscr V_{\ell+1}}_{\psi_Q(t)}\left(\mathsf F_\ell Q\right)
\quad\text{for every }t\in\mathcal T_Q,
\label{eq:hist-oriented-flow-semiconjugacy}
\end{equation}
where $\bar{\mathcal T}^{\max}_{\mathsf F_\ell Q}=(\bar t^-,\bar t^+)$ is the maximal
interval of the coarse flow through $\mathsf F_\ell Q$.  With $\mathcal T_Q=(t^-,t^+)$
and endpoints in the extended real line,
\begin{equation}
\mathcal T'_Q
=\left(
\lim_{t\downarrow t^-}\psi_Q(t),
\lim_{t\uparrow t^+}\psi_Q(t)
\right),
\label{eq:hist-semiconjugacy-endpoint-image}
\end{equation}
and the whole maximal coarse orbit is traversed if and only if the two limits equal
$\bar t^-$ and $\bar t^+$; when the fine orbit is complete in $\mathcal U$ and the coarse
maximal interval is $\R$, this means
$\int_0^{+\infty}\mathsf a_\ell(\operatorname{Fl}^{\mathscr V_\ell}_{t}Q)dt
=\int_{-\infty}^{0}\mathsf a_\ell(\operatorname{Fl}^{\mathscr V_\ell}_{t}Q)dt=+\infty$.
A positive lower bound on $\mathsf a_\ell$ suffices and is not necessary: on $\R$ with
$\mathscr V_\ell=\mathscr V_{\ell+1}=\partial_z$ and
$\mathsf F_\ell=\operatorname{arsinh}$, $\mathsf a_\ell=(1+z^2)^{-1/2}$ has infimum zero
while $\psi_0=\operatorname{arsinh}$ maps $\R$ onto $\R$; with $\mathsf F_\ell=\arctan$,
$\mathsf a_\ell=(1+z^2)^{-1}$ is positive while $\mathcal T'_0=(-\pi/2,\pi/2)$ although
both flows are complete.  Conversely, differentiating a flow relation of the form
\eqref{eq:hist-oriented-flow-semiconjugacy} at $t=0$ gives
\eqref{eq:hist-oriented-vector-field-semiconjugacy}.  \status{ESTABLISHED}

The proof is in \cref{sec:proofs-fiber-geometry}.

\propositionheading{Noncollapse is a separate hypothesis}{prop:hist-noncollapse}
Condition \eqref{eq:hist-oriented-vector-field-semiconjugacy} alone permits total
collapse: with $\mathsf F_\ell$ constant and $\mathscr V_{\ell+1}\equiv0$ it holds for
every positive $\mathsf a_\ell$, while every fine orbit maps to a point and every coarse
duration is zero.  A nonconstant shared oriented history additionally requires
$T_Q\mathsf F_\ell\mathscr V_\ell(Q)\ne0$ on every nontrivial subarc, equivalently
$\mathscr V_{\ell+1}(\mathsf F_\ell Q)\ne0$ there, and full traversal requires the
endpoint equalities of \eqref{eq:hist-semiconjugacy-endpoint-image}; neither follows from
\eqref{eq:hist-oriented-vector-field-semiconjugacy}.  \status{ESTABLISHED}

\noindent\emph{Proof.}
Both sides of \eqref{eq:hist-oriented-vector-field-semiconjugacy} vanish in the example,
and the arctan example of \cref{thm:hist-oriented-semiconjugacy} satisfies it without
full traversal.  $\square$

\theoremheading{The semiconjugacy rate cancels from geometric length}{thm:hist-duration-relation}
Under \eqref{eq:hist-oriented-vector-field-semiconjugacy}, let $Q(t)$ be a fine orbit.
The image speed is
\begin{equation}
\left\lVert\tfrac{d}{dt}\mathsf F_\ell Q(t)\right\rVert_{\mathsf G_{\ell+1}}
=\mathsf a_\ell\left(Q(t)\right)
\left\lVert\mathscr V_{\ell+1}\left(\mathsf F_\ell Q(t)\right)\right\rVert_{\mathsf G_{\ell+1}},
\label{eq:hist-image-speed}
\end{equation}
and for $t_0<t_1$ in $\mathcal T_Q$ the accumulated coarse duration is
\begin{equation}
\mathfrak t^{(\ell+1)}(t_1)-\mathfrak t^{(\ell+1)}(t_0)
=\int_{\psi_Q(t_0)}^{\psi_Q(t_1)}
\left\lVert\mathscr V_{\ell+1}\left(\operatorname{Fl}^{\mathscr V_{\ell+1}}_{t'}(\mathsf F_\ell Q)\right)
\right\rVert_{\mathsf G_{\ell+1}}dt',
\label{eq:hist-coarse-duration}
\end{equation}
the intrinsic $\mathsf G_{\ell+1}$-length of the coarse orbit arc, which does not depend
on $\mathsf a_\ell$.  \status{ESTABLISHED}

\noindent\emph{Proof.}
Equation \eqref{eq:hist-image-speed} is \eqref{eq:hist-oriented-vector-field-semiconjugacy}
with positive homogeneity of the norm and $\mathsf a_\ell>0$.  Substituting it and
changing variables $t'=\psi_Q(t)$, whose Jacobian is
$\mathsf a_\ell(\operatorname{Fl}^{\mathscr V_\ell}_tQ)$, cancels the factor.  $\square$

\remarkheading{Exact duration comparison}{thm:hist-duration-criterion}
Under \eqref{eq:hist-oriented-vector-field-semiconjugacy}, along one orbit and from a
common origin, the coarse and fine durations agree identically if and only if
$\lVert T\mathsf F_\ell\mathscr V_\ell\rVert_{\mathsf G_{\ell+1}}=\lVert\mathscr
V_\ell\rVert_{\mathsf G_\ell}$ almost everywhere along the orbit, they agree up to a
factor $\mathrm c>0$ if and only if the same holds with that factor, and the coarse
duration has no larger increments on every subinterval if and only if
\begin{equation}
\left\lVert T\mathsf F_\ell\mathscr V_\ell\right\rVert_{\mathsf G_{\ell+1}}
\leq\left\lVert\mathscr V_\ell\right\rVert_{\mathsf G_\ell}
\quad\text{along the orbit},
\label{eq:hist-duration-criterion}
\end{equation}
because both durations are integrals of continuous nonnegative densities from a common
origin, and two such integrals agree for every upper limit exactly when the densities
agree almost everywhere.  The condition is scalar and concerns the single direction
$\mathscr V_\ell$; the Loewner condition $\mathsf F_\ell^*\mathsf G_{\ell+1}\preceq\mathsf
G_\ell$ is sufficient and stronger, and \cref{thm:hist-pointwise-contraction-lift}
supplies it for channelwise Markov maps.  Fiberwise Fisher contraction alone gives
neither, for two independent reasons.  The configuration metrics are separate data: on
$\R$ with $\mathscr V_\ell=\mathscr V_{\ell+1}=\partial_z$ and $\mathsf F_\ell$ the
identity, the semiconjugacy holds with $\mathsf a_\ell\equiv1$, and the metrics $dz^2$
and $4dz^2$ double the duration, whether realized with the identity fiber map and unequal
channel weights $1$ and $4$ (zero Fisher defect) or with fibers $\mathcal N(\mu,1)$ and
$\mathcal N(\mu,\tfrac14)$ and the identity parameter map (Fisher defect $1-4=-3$, so no
Markov morphism).  And contraction compares the image of the fine orbit, not a recomputed
coarse orbit: on $\R^2$ with $\mathsf F_\ell$ the identity and the Euclidean metric at
both depths, $\Fenergy_\ell=\tfrac12(z_1^2+z_2^2)$ and
$\Fenergy_{\ell+1}=\tfrac12(z_1^2+4z_2^2)$ give the fine history $(e^{-t},e^{-t})$ of
duration $\sqrt2$ and the recomputed coarse history $(e^{-t},e^{-4t})$ of duration
$\int_0^1\sqrt{1+16z^6} dz>\sqrt2$, and
\eqref{eq:hist-oriented-vector-field-semiconjugacy} fails because
$(z_1,z_2)=\mathsf a(z_1,4z_2)$ forces $\mathsf a=1$ and $z_2=0$.  \status{ESTABLISHED}

\corollaryheading{When fiberwise contraction does lift}{thm:hist-pointwise-contraction-lift}
Let $\mathcal Q_\ell$ and $\mathcal Q_{\ell+1}$ carry the channel-weighted metrics
$\mathsf G_\ell=\sum_xw_xg^F_x$ and $\mathsf G_{\ell+1}=\sum_xw_x\bar g^F_x$ with the same
positive weights, and let $\mathsf F_\ell$ act channelwise by Markov coarse morphisms,
$(\mathsf F_\ell Q)_x=\mathscr K_x(Q_x)$.  Then for every $Z\in T_Q\mathcal Q_\ell$,
\begin{equation}
\lVert Z\rVert^2_{\mathsf G_\ell}
-\lVert T\mathsf F_\ell Z\rVert^2_{\mathsf G_{\ell+1}}
=\sum_xw_x\Delta_F^{\mathscr K_x}(Z_x,Z_x)
\geq0,
\label{eq:hist-pointwise-contraction-lift}
\end{equation}
so the image of every $C^1$ fine curve has no greater speed and no greater duration on
every subinterval, and \eqref{eq:hist-duration-criterion} holds along every orbit of an
oriented semiconjugacy. On a subinterval of finite fine duration, equality holds exactly
when, for every channel and almost every parameter value, the fine score in the
direction $Z_x$ is measurable with respect to the coarse output.
\status{ESTABLISHED}

\noindent\emph{Proof.}
Differentiate the channelwise action, insert it into $\mathsf G_{\ell+1}$, and subtract
from $\mathsf G_\ell$; the difference is the weighted sum of the vertical defects, which
are positive semidefinite with the stated equality case by
\cref{thm:pb-pullback-fisher-defect}.  Monotonicity of the square root and of the
integral gives the speed and duration statements.  $\square$

\section{Semiconjugacy on symmetry quotients}
\label{sec:hist-natural-gradient-semiconjugacy}

Throughout this section $\mathscr V_\ell$ and $\mathscr V_{\ell+1}$ are natural-gradient
fields of $C^2$ objectives for positive-definite metrics.

\definitionheading{Functional compatibility}{hyp:hist-functional-compatibility}
A coarse configuration map $\mathsf F_\ell$ is functionally compatible on an open set
$\mathcal U\subseteq\mathcal Q_\ell$ when there is a $C^1$ function $\chi_\ell$ with
$\chi_\ell'>0$ such that $\Fenergy_\ell=\chi_\ell\circ\Fenergy_{\ell+1}\circ\mathsf
F_\ell$ on $\mathcal U$.

\propositionheading{Semiconjugacy on symmetry quotients}{prop:hist-natural-gradient-sufficiency}
Equal objectives do not intertwine natural gradients: with $\mathsf F_\ell$ the identity
of $\R^2$, $\Fenergy_\ell=\Fenergy_{\ell+1}=\tfrac12(z_1^2+2z_2^2)$, and metrics
$\operatorname{diag}(1,1)$ and $\operatorname{diag}(1,\mathrm c)$ with $\mathrm c\ne1$,
\eqref{eq:hist-oriented-vector-field-semiconjugacy} fails wherever $z_1z_2\ne0$.
Functional compatibility on $\mathcal U$ forces $d\Fenergy_\ell$ to vanish on
$\ker T\mathsf F_\ell$ at every point of $\mathcal U$, so when $\mathsf F_\ell$ is a
submersion, $\Fenergy_\ell$ is constant on each connected piece of each fiber of
$\mathsf F_\ell$ in $\mathcal U$.  If $\mathsf F_\ell$ is a surjective submersion,
functionally compatible on $\mathcal U$, and horizontally conformal with dilation
$\mathsf d_\ell>0$, meaning that
$\mathsf G_{\ell+1}(T\mathsf F_\ell Z_1,T\mathsf F_\ell Z_2)=\mathsf d_\ell^2\mathsf
G_\ell(Z_1,Z_2)$ for $Z_1,Z_2\in\operatorname{Hor}_Q=(\ker T_Q\mathsf F_\ell)^\perp$, then
\eqref{eq:hist-oriented-vector-field-semiconjugacy} holds on $\mathcal U$ with
\begin{equation}
\mathsf a_\ell(Q)
=\chi_\ell'\left(\Fenergy_{\ell+1}(\mathsf F_\ell Q)\right)
\mathsf d_\ell(Q)^2>0 .
\label{eq:hist-conformal-semiconjugacy-factor}
\end{equation}
In particular the quotient map of a free, proper, isometric action that preserves
$\Fenergy_\ell$, with the quotient metric and descended objective of
\cref{prop:hist-invariant-history-horizontal}, is an oriented semiconjugacy with
$\mathsf a_\ell\equiv1$; and at $c_*$ the identity map from the mean coordinate of the
fiber $\{\mathcal N(\mu,1)\}$ to that of $\{\mathcal N(\mu,2)\}$, realized by the kernel
$y\mapsto\mathcal N(y,1)$, is horizontally conformal with $\mathsf d_\ell^2=\tfrac12$ and
gives $\mathsf a_\ell=\tfrac12$ for $\Fenergy_\ell=\Fenergy_{\ell+1}=\mu^2/2$, an
injective case that does not exercise the fiber condition.  The route does not reach
maps that discard directions on which the fine objective depends: if
$\Fenergy_{\ell+1}(Q')=\inf\{\Fenergy_\ell(Q):\mathsf F_\ell Q=Q'\}$, $\mathsf F_\ell$ is
a submersion with positive-dimensional fibers, and each fiber contains at most one
minimizer of $\Fenergy_\ell$, then compatibility with $\chi_\ell=\operatorname{id}$ holds
on no open set; for $\mathsf F_\ell(f,s)=s$ and
$\Fenergy_\ell(f,s)=\tfrac12(f-s)^2+\tfrac12s^2$ it holds for no $\chi_\ell$ on any open
set.  \status{ESTABLISHED}

The proof is in \cref{sec:proofs-fiber-geometry}.

For maps that discard directions on which the fine objective depends, the exact route
runs through the reduced slow flow, which on the set of conditional minimizers is
conjugate to the coarse natural-gradient flow (\cref{prop:obs-rate-separated-flow},
\eqref{eq:obs-reduced-semiconjugacy}); whether the fine or rate-separated flow follows
it is \cref{open:obs-slow-manifold}.

\section{Coarse-graining depth, orbit parameter, and duration}
\label{sec:hist-rg-depth-not-time}

A tower of coarse configurations $Q^{(\ell+1)}=\mathsf F_\ell(Q^{(\ell)})$,
$\ell=0,1,2,\ldots$, is indexed by the coarse-graining depth $\ell$, which is not a
parameter on a history.  A history can be present at every depth, as $Q^{(\ell)}(t)$
with accumulated duration $\mathfrak t^{(\ell)}(t)$; when
$Q^{(\ell+1)}(t)=\mathsf F_\ell(Q^{(\ell)}(t))$ and $\mathsf F_\ell$ acts by Markov coarse
morphisms as in \cref{thm:hist-pointwise-contraction-lift}, the record-clock comparison
applies between adjacent depths, while for a general oriented semiconjugacy the exact
comparison is \cref{thm:hist-duration-criterion}, and independently recomputed
histories also need the noncollapse conditions of \cref{prop:hist-noncollapse}.  The
Markov property of a configuration arrow can fail even for natural meta-agent maps:
no parameter-independent Markov kernel sends every product of two normal laws to
their joint mean-and-variance barycenter, as Supplement~G proves.

\propositionheading{Typed non-identification of depth, orbit parameter, and duration}{prop:hist-coordinate-independence}
Coarse-graining depth $\ell$, an orbit parameter $t$, and Fisher duration
$\mathfrak t^{(\ell)}$ are distinct: $\ell$ is a discrete index, $t$ a chosen oriented
parameter on a selected orbit, and $\mathfrak t^{(\ell)}$ the arc length for a chosen
metric from a chosen origin.  Once the metric, origin, orientation, and parameterized path
are chosen, $\mathfrak t^{(\ell)}$ is a function of $t$, but it is not determined by the
unparameterized orbit or by $\ell$: changing the metric rescales it and changing the
origin shifts it.  Strict monotonicity of $\mathfrak t^{(\ell)}$ does not make it a
regular coordinate: for $\mathsf G=dz^2$ on $\R$ and the path $z(t)=t^2$ on $[0,1]$,
$\mathfrak t(t)=t^2$, whose inverse is not differentiable at the origin, and a zero-speed
subinterval destroys strict monotonicity.  \status{ESTABLISHED}

\noindent\emph{Proof.}
Only the metric, origin, orientation, and parameterized path produce a duration
function.  Metric dependence is homogeneity of the norm in \eqref{eq:hist-fisher-speed},
and the example is $\int_0^t2t''dt''=t^2$ with the square root not differentiable at
zero.  $\square$

\section*{Physicist's summary}

The Fisher metric on an agent's belief and model law fibers is a gauge-invariant ruler and
the second-order term of the relative entropy, so a loop holonomy near the identity
appears as a quadratic form on the Lie algebra that vanishes on the isotropy algebra of
the law.  A Markov channel loses exactly the conditional variance of the score, so a
meta-agent that is a channel image of its constituents' history has an information clock
no faster than theirs, by a ratio that is open.  A chosen metric turns the free energy
into an orbit, Fisher arc length removes its parameterization, and a belief-model rate
ratio, unlike a scalar mobility, can bend the orbit; an orthogonal clock needs the
gradient norm to be locally a function of the free energy, a unit-speed clock only a
global cross-section.  By the parity witness, joint laws, Fisher metrics, and free
energies are not functions of marginals.  A recomputed coarse flow is the image of the
fine flow only under an oriented semiconjugacy, which holds with rate one on symmetry
quotients; information-discarding maps are compared through the reduced slow flow of the
next chapter.
